\chapter{Chapter07 — CSP}
%%% generated by the blueprint pipeline from FABL.Chapter07.CSP
\section{Overview}

Declarations of the Lean module \texttt{FABL.Chapter07.CSP}.

\section{Declarations}

\begin{definition}[Constraint satisfaction template]
\lean{FABL.CSPTemplate}
\statementsource{Analysis-of-Boolean-Functions-by-Ryan-ODonnell}{pdf-ffbe9c7b3306-p184-b004}
\label{FABL.CSPTemplate}
\leanok
A constraint satisfaction template over a domain $D$ consists of a finite set of
predicate names, an arity for each predicate, and a Boolean evaluation function for each
predicate on ordered tuples of domain elements of that arity; the value true means the
constraint is satisfied.
\end{definition}

\begin{definition}[Maximum predicate arity]
\lean{FABL.CSPTemplate.maxArity}
\statementsource{Analysis-of-Boolean-Functions-by-Ryan-ODonnell}{pdf-ffbe9c7b3306-p184-b004}
\label{FABL.CSPTemplate.maxArity}
\leanok
For a constraint satisfaction template with a finite set of predicates, its maximum
arity is the supremum of the arities of all its predicates, taking the supremum over the
entire predicate set.
\end{definition}

\begin{definition}[Unsatisfied predicate cost]
\lean{FABL.CSPTemplate.unsatisfiedCost}
\label{FABL.CSPTemplate.unsatisfiedCost}
\leanok
Let $\Psi$ be a constraint satisfaction template with domain $D$, let $p$ be one of its
predicates, and let $x$ be an input tuple of the predicate's arity. The unsatisfied cost
is $0$ if $p$ accepts $x$, and $1$ otherwise.
\end{definition}

\begin{definition}[Valued CSP of a template]
\lean{FABL.CSPTemplate.toValuedCSP}
\label{FABL.CSPTemplate.toValuedCSP}
\leanok
\uses{FABL.CSPTemplate.unsatisfiedCost}
Let $\Psi$ be a constraint satisfaction template over domain $D$. Its associated valued
CSP consists of all functions obtained from a predicate $p$ of $\Psi$ by pairing that
predicate's arity with its unsatisfied cost function.
\end{definition}

\begin{theorem}[Maximum arity for a constant-arity family]
\lean{FABL.CSPTemplate.maxArity_eq_of_forall_arity_eq}
\label{FABL.CSPTemplate.maxArity_eq_of_forall_arity_eq}
\leanok
\uses{FABL.CSPTemplate.maxArity}
\difficulty{4}
Let a nonempty finite family of predicates on a domain be given, and suppose every
predicate has arity $r$, where $r$ is a natural number. Then the maximum arity in the
family is $r$.
\end{theorem}

\begin{definition}[Constraint satisfaction constraint]
\lean{FABL.CSPConstraint}
\statementsource{Analysis-of-Boolean-Functions-by-Ryan-ODonnell}{pdf-ffbe9c7b3306-p185-b004}
\label{FABL.CSPConstraint}
\leanok
\uses{FABL.CSPTemplate}
Given a constraint satisfaction template and a set of variables, a constraint consists
of one allowed predicate together with an ordered tuple of distinct variables whose
length is the arity of that predicate.
\end{definition}

\begin{definition}[Satisfied constraint]
\lean{FABL.CSPConstraint.Satisfied}
\statementsource{Analysis-of-Boolean-Functions-by-Ryan-ODonnell}{pdf-ffbe9c7b3306-p186-b001}
\label{FABL.CSPConstraint.Satisfied}
\leanok
\uses{FABL.CSPConstraint}
Given a constraint, a constraint satisfaction template, and an assignment of domain
values to the variables, the constraint is satisfied exactly when the template's
predicate evaluates to true on the tuple of assigned values in the constraint's ordered
scope.
\end{definition}

\begin{definition}[Constraint satisfaction indicator]
\lean{FABL.CSPConstraint.satisfiedIndicator}
\label{FABL.CSPConstraint.satisfiedIndicator}
\leanok
\uses{FABL.CSPConstraint}
Given a constraint and an assignment of values to its variables, the satisfaction
indicator is $1$ when the constraint’s predicate evaluates to true on the values
assigned to its ordered scope, and $0$ otherwise.
\end{definition}

\begin{definition}[Constraint violation indicator]
\lean{FABL.CSPConstraint.unsatisfiedIndicator}
\label{FABL.CSPConstraint.unsatisfiedIndicator}
\leanok
\uses{FABL.CSPConstraint}
Given a constraint and an assignment of values to its variables, the violation indicator
is $0$ when the constraint’s predicate evaluates to true on the values assigned to its
ordered scope, and $1$ otherwise.
\end{definition}

\begin{definition}[Valued term associated to a constraint]
\lean{FABL.CSPConstraint.toValuedTerm}
\label{FABL.CSPConstraint.toValuedTerm}
\leanok
\uses{FABL.CSPConstraint, FABL.CSPTemplate.toValuedCSP, FABL.CSPTemplate.unsatisfiedCost}
A constraint determines a valued term whose arity is the arity of its predicate, whose
cost function is that predicate’s unsatisfied-cost function, and whose variables are the
constraint’s ordered scope.
\end{definition}

\begin{theorem}[Constraint cost on an assignment]
\lean{FABL.CSPConstraint.evalSolution_toValuedTerm}
\label{FABL.CSPConstraint.evalSolution_toValuedTerm}
\leanok
\uses{FABL.CSPConstraint, FABL.CSPConstraint.toValuedTerm, FABL.CSPConstraint.unsatisfiedIndicator, FABL.CSPTemplate.unsatisfiedCost}
\difficulty{1}
Let a constraint template specify Boolean-valued predicates, let $c$ be one of its
constraints, and let $\sigma$ assign a domain value to every variable. The cost of the
valued term associated to $c$, evaluated at $\sigma$, equals $0$ if $c$ is satisfied and
$1$ if it is not.
\end{theorem}

\begin{theorem}[Satisfied and unsatisfied indicators sum to one]
\lean{FABL.CSPConstraint.satisfiedIndicator_add_unsatisfiedIndicator}
\label{FABL.CSPConstraint.satisfiedIndicator_add_unsatisfiedIndicator}
\leanok
\uses{FABL.CSPConstraint, FABL.CSPConstraint.satisfiedIndicator, FABL.CSPConstraint.unsatisfiedIndicator}
\difficulty{4}
Let a constraint template on a set $D$ specify a finite family of predicates with
arities and Boolean evaluations, and let $c$ be one of its constraints on a set of
variables. For every assignment of values in $D$ to the variables, the constraint’s
satisfied indicator plus its unsatisfied indicator equals $1$.
\end{theorem}

\begin{definition}[Constraint satisfaction instance]
\lean{FABL.CSPInstance}
\statementsource{Analysis-of-Boolean-Functions-by-Ryan-ODonnell}{pdf-ffbe9c7b3306-p185-b004}
\label{FABL.CSPInstance}
\leanok
\uses{FABL.CSPConstraint, FABL.CSPTemplate}
Let $D$ be a domain, let a finite family of predicates on $D$ be specified with an arity
and a Boolean evaluation rule for each predicate, and let $V$ be a set of variables. A
constraint instance consists of a nonempty multiset of constraints, each applying one of
the specified predicates to an ordered tuple of distinct variables, such that every
variable in $V$ occurs in the scope of at least one constraint.
\end{definition}

\begin{definition}[Valued instance associated to a CSP instance]
\lean{FABL.CSPInstance.toValuedInstance}
\label{FABL.CSPInstance.toValuedInstance}
\leanok
\uses{FABL.CSPConstraint.toValuedTerm, FABL.CSPInstance, FABL.CSPTemplate.toValuedCSP}
A CSP instance determines a valued instance by mapping each constraint in its multiset
to the valued term associated to that constraint, preserving multiplicities.
\end{definition}

\begin{definition}[Number of satisfied constraints]
\lean{FABL.CSPInstance.satisfiedCount}
\label{FABL.CSPInstance.satisfiedCount}
\leanok
\uses{FABL.CSPConstraint.satisfiedIndicator, FABL.CSPInstance}
For a CSP instance and an assignment of values to its variables, the satisfied count is
the sum, with multiplicities, of the satisfaction indicators of its constraints.
\end{definition}

\begin{definition}[Total unsatisfied cost]
\lean{FABL.CSPInstance.unsatisfiedCost}
\label{FABL.CSPInstance.unsatisfiedCost}
\leanok
\uses{FABL.CSPInstance, FABL.CSPInstance.toValuedInstance}
For a CSP instance and an assignment of values to its variables, the unsatisfied cost is
the value of the associated valued instance on that assignment.
\end{definition}

\begin{theorem}[Unsatisfied Cost as an Indicator Sum]
\lean{FABL.CSPInstance.unsatisfiedCost_eq_sum}
\label{FABL.CSPInstance.unsatisfiedCost_eq_sum}
\leanok
\uses{FABL.CSPConstraint.unsatisfiedIndicator, FABL.CSPInstance, FABL.CSPInstance.toValuedInstance, FABL.CSPInstance.unsatisfiedCost}
\difficulty{1}
For any constraint satisfaction instance and any assignment of domain values to its
variables, the total unsatisfied cost equals the sum of the unsatisfaction indicators
over its constraints, counting multiplicities.
\end{theorem}

\begin{theorem}[Constraint Indicator Sum]
\lean{FABL.CSPInstance.sum_indicators_eq_card}
\label{FABL.CSPInstance.sum_indicators_eq_card}
\leanok
\uses{FABL.CSPConstraint, FABL.CSPConstraint.satisfiedIndicator, FABL.CSPConstraint.satisfiedIndicator_add_unsatisfiedIndicator, FABL.CSPConstraint.unsatisfiedIndicator}
\difficulty{4}
Let $D$ be a domain, let $[?]$
\end{theorem}

\begin{theorem}[Satisfied Count and Unsatisfied Cost]
\lean{FABL.CSPInstance.satisfiedCount_add_unsatisfiedCost}
\label{FABL.CSPInstance.satisfiedCount_add_unsatisfiedCost}
\leanok
\uses{FABL.CSPInstance, FABL.CSPInstance.satisfiedCount, FABL.CSPInstance.sum_indicators_eq_card, FABL.CSPInstance.unsatisfiedCost, FABL.CSPInstance.unsatisfiedCost_eq_sum}
\difficulty{3}
For any constraint satisfaction instance and any assignment of domain values to its
variables, the number of satisfied constraints plus the total unsatisfied cost equals
the number of constraints, counting multiplicities.
\end{theorem}

\begin{definition}[Fraction of satisfied constraints]
\lean{FABL.CSPInstance.value}
\statementsource{Analysis-of-Boolean-Functions-by-Ryan-ODonnell}{
  pdf-ffbe9c7b3306-p185-b006,
  pdf-ffbe9c7b3306-p186-b001
}
\label{FABL.CSPInstance.value}
\leanok
\uses{FABL.CSPInstance, FABL.CSPInstance.satisfiedCount}
For a constraint instance and an assignment of domain elements to its variables, define
its value to be the number of satisfied constraints divided by the total number of
constraints, counting multiplicities.
\end{definition}

\begin{theorem}[Value from Unsatisfied Cost]
\lean{FABL.CSPInstance.value_eq_one_sub_unsatisfiedCost_div}
\label{FABL.CSPInstance.value_eq_one_sub_unsatisfiedCost_div}
\leanok
\uses{FABL.CSPInstance, FABL.CSPInstance.satisfiedCount, FABL.CSPInstance.satisfiedCount_add_unsatisfiedCost, FABL.CSPInstance.unsatisfiedCost, FABL.CSPInstance.value}
\difficulty{4}
For any constraint satisfaction instance and any assignment of domain values to its
variables, the fraction of satisfied constraints equals one minus the total unsatisfied
cost divided by the number of constraints.
\end{theorem}

\begin{theorem}[Nonnegativity of the constraint satisfaction value]
\lean{FABL.CSPInstance.value_nonneg}
\label{FABL.CSPInstance.value_nonneg}
\leanok
\uses{FABL.CSPInstance, FABL.CSPInstance.value}
\difficulty{1}
Let $D$ be a domain and let a finite family of predicates on $D$ be given, each with a
specified arity and Boolean evaluation rule. Consider an instance with a nonempty
multiset of constraints, in which every declared variable occurs in some constraint, and
any assignment of domain elements to the variables. The fraction of constraints
satisfied by the assignment is nonnegative.
\end{theorem}

\begin{theorem}[Bound on the Satisfied Fraction]
\lean{FABL.CSPInstance.value_le_one}
\label{FABL.CSPInstance.value_le_one}
\leanok
\uses{FABL.CSPInstance, FABL.CSPInstance.value, FABL.CSPInstance.value_eq_one_sub_unsatisfiedCost_div}
\difficulty{3}
For any constraint satisfaction instance and any assignment of domain values to its
variables, the fraction of satisfied constraints is at most $1$.
\end{theorem}

\begin{definition}[Full satisfaction]
\lean{FABL.CSPInstance.FullySatisfied}
\statementsource{Analysis-of-Boolean-Functions-by-Ryan-ODonnell}{pdf-ffbe9c7b3306-p186-b001}
\label{FABL.CSPInstance.FullySatisfied}
\leanok
\uses{FABL.CSPConstraint.Satisfied, FABL.CSPInstance}
For a constraint instance and an assignment of domain elements to its variables, the
instance is fully satisfied if every constraint in its multiset is satisfied by that
assignment.
\end{definition}

\begin{theorem}[Satisfaction and zero unsatisfied count]
\lean{FABL.CSPInstance.forall_satisfied_iff_sum_unsatisfied_eq_zero}
\label{FABL.CSPInstance.forall_satisfied_iff_sum_unsatisfied_eq_zero}
\leanok
\uses{FABL.CSPConstraint, FABL.CSPConstraint.Satisfied, FABL.CSPConstraint.unsatisfiedIndicator}
\difficulty{4}
For any multiset of constraints from a fixed finite-predicate constraint template and
any assignment of labels to variables, every constraint in the multiset is satisfied if
and only if the sum of the unsatisfied indicators, counted with multiplicity, is zero.
\end{theorem}

\begin{theorem}[Full satisfaction and zero cost]
\lean{FABL.CSPInstance.fullySatisfied_iff_unsatisfiedCost_eq_zero}
\label{FABL.CSPInstance.fullySatisfied_iff_unsatisfiedCost_eq_zero}
\leanok
\uses{FABL.CSPInstance, FABL.CSPInstance.FullySatisfied, FABL.CSPInstance.forall_satisfied_iff_sum_unsatisfied_eq_zero, FABL.CSPInstance.unsatisfiedCost, FABL.CSPInstance.unsatisfiedCost_eq_sum}
\difficulty{1}
Let a constraint satisfaction instance have a nonempty multiset of constraints in which
every variable occurs, and let an assignment of labels be given. The assignment
satisfies every constraint if and only if its total unsatisfied cost, counted with
constraint multiplicities, is zero.
\end{theorem}

\begin{theorem}[Value eq one iff fully satisfied]
\lean{FABL.CSPInstance.value_eq_one_iff_fullySatisfied}
\statementsource{Analysis-of-Boolean-Functions-by-Ryan-ODonnell}{pdf-ffbe9c7b3306-p186-b001}
\label{FABL.CSPInstance.value_eq_one_iff_fullySatisfied}
\leanok
\uses{FABL.CSPInstance, FABL.CSPInstance.FullySatisfied, FABL.CSPInstance.fullySatisfied_iff_unsatisfiedCost_eq_zero, FABL.CSPInstance.value, FABL.CSPInstance.value_eq_one_sub_unsatisfiedCost_div}
\difficulty{4}
% dataset-pipeline: stub -- prose pending build_informal_statements.py
Statement pending regeneration.
\end{theorem}

\begin{definition}[Optimal assignment]
\lean{FABL.CSPInstance.optimalAssignment}
\statementsource{Analysis-of-Boolean-Functions-by-Ryan-ODonnell}{pdf-ffbe9c7b3306-p186-b001}
\label{FABL.CSPInstance.optimalAssignment}
\leanok
\uses{FABL.CSPInstance, FABL.CSPInstance.value}
For a constraint instance with a finite, nonempty domain and a finite variable set, an
optimal assignment is an assignment attaining the maximum possible value of the
instance.
\end{definition}

\begin{definition}[Optimal value]
\lean{FABL.CSPInstance.optimum}
\statementsource{Analysis-of-Boolean-Functions-by-Ryan-ODonnell}{pdf-ffbe9c7b3306-p186-b001}
\label{FABL.CSPInstance.optimum}
\leanok
\uses{FABL.CSPInstance, FABL.CSPInstance.optimalAssignment, FABL.CSPInstance.value}
For a constraint instance with a finite, nonempty domain and a finite variable set, its
optimum is the value attained by an optimal assignment.
\end{definition}

\begin{theorem}[Assignment Value Bounded by the Optimum]
\lean{FABL.CSPInstance.value_le_optimum}
\label{FABL.CSPInstance.value_le_optimum}
\leanok
\uses{FABL.CSPInstance, FABL.CSPInstance.optimum, FABL.CSPInstance.value}
\difficulty{1}
Let a constraint satisfaction instance have a finite, nonempty domain and a finite
variable set. For every assignment of domain values to its variables, its fraction of
satisfied constraints is at most the maximum such fraction over all assignments.
\end{theorem}

\begin{theorem}[Attainment of the optimum]
\lean{FABL.CSPInstance.exists_value_eq_optimum}
\label{FABL.CSPInstance.exists_value_eq_optimum}
\leanok
\uses{FABL.CSPInstance, FABL.CSPInstance.optimalAssignment, FABL.CSPInstance.optimum, FABL.CSPInstance.value}
\difficulty{1}
Let $D$ be a finite nonempty label set, $V$ a finite variable set, and consider a
constraint satisfaction instance on $V$ with a nonempty multiset of constraints in which
every variable occurs. There is an assignment of labels in $D$ to the variables whose
fraction of satisfied constraints equals the instance's optimum value.
\end{theorem}

\begin{theorem}[Nonnegativity of the optimum]
\lean{FABL.CSPInstance.optimum_nonneg}
\label{FABL.CSPInstance.optimum_nonneg}
\leanok
\uses{FABL.CSPInstance, FABL.CSPInstance.optimalAssignment, FABL.CSPInstance.optimum, FABL.CSPInstance.value_nonneg}
\difficulty{1}
Let a finite-domain constraint satisfaction problem have a nonempty multiset of
constraints, with every declared variable occurring in some constraint scope. If both
the domain and the variable set are finite and the domain is nonempty, then the optimum
value is nonnegative.
\end{theorem}

\begin{theorem}[Upper bound on the optimum]
\lean{FABL.CSPInstance.optimum_le_one}
\label{FABL.CSPInstance.optimum_le_one}
\leanok
\uses{FABL.CSPInstance, FABL.CSPInstance.optimalAssignment, FABL.CSPInstance.optimum, FABL.CSPInstance.value_le_one}
\difficulty{1}
Let a finite-domain constraint satisfaction problem have a nonempty multiset of
constraints, with every declared variable occurring in some constraint scope. If both
the domain and the variable set are finite and the domain is nonempty, then the optimum
value is at most $1$.
\end{theorem}

\begin{definition}[Satisfiability by optimum]
\lean{FABL.CSPInstance.Satisfiable}
\statementsource{Analysis-of-Boolean-Functions-by-Ryan-ODonnell}{pdf-ffbe9c7b3306-p186-b001}
\label{FABL.CSPInstance.Satisfiable}
\leanok
\uses{FABL.CSPInstance, FABL.CSPInstance.optimum}
For a constraint instance with a finite, nonempty domain and a finite variable set, the
instance is satisfiable precisely when its optimum equals $1$.
\end{definition}

\begin{theorem}[Satisfiability and Full Satisfaction]
\lean{FABL.CSPInstance.satisfiable_iff_exists_fullySatisfied}
\label{FABL.CSPInstance.satisfiable_iff_exists_fullySatisfied}
\leanok
\uses{FABL.CSPInstance, FABL.CSPInstance.FullySatisfied, FABL.CSPInstance.Satisfiable, FABL.CSPInstance.optimalAssignment, FABL.CSPInstance.optimum_le_one, FABL.CSPInstance.value_eq_one_iff_fullySatisfied, FABL.CSPInstance.value_le_optimum}
\difficulty{4}
Let a constraint satisfaction instance have a finite, nonempty domain and a finite
variable set. The instance is satisfiable, meaning its maximum fraction of satisfied
constraints is $1$, if and only if there is an assignment satisfying every constraint.
\end{theorem}

\begin{definition}[Uniform string tester]
\lean{FABL.UniformStringTester}
\statementsource{Analysis-of-Boolean-Functions-by-Ryan-ODonnell}{pdf-ffbe9c7b3306-p187-b001}
\label{FABL.UniformStringTester}
\leanok
\uses{FABL.CSPTemplate}
Let $D$ be a domain, let a finite family of predicates on $D$ be specified with an arity
for each predicate, and let $V$ be a set of coordinates. A uniform string tester
consists of a finite, nonempty seed space, a choice of predicate for each seed, and an
ordered tuple of distinct queried coordinates of the required arity for each seed, with
every coordinate in $V$ queried for at least one seed.
\end{definition}

\begin{definition}[Constraint selected by a seed]
\lean{FABL.UniformStringTester.constraint}
\statementsource{Analysis-of-Boolean-Functions-by-Ryan-ODonnell}{pdf-ffbe9c7b3306-p187-b001}
\label{FABL.UniformStringTester.constraint}
\leanok
\uses{FABL.CSPConstraint, FABL.UniformStringTester}
Given a uniform string tester and a seed, the associated constraint uses the predicate
selected by that seed and the ordered tuple of coordinates queried by that seed.
\end{definition}

\begin{definition}[CSP instance induced by a tester]
\lean{FABL.UniformStringTester.toCSPInstance}
\statementsource{Analysis-of-Boolean-Functions-by-Ryan-ODonnell}{pdf-ffbe9c7b3306-p187-b001}
\label{FABL.UniformStringTester.toCSPInstance}
\leanok
\uses{FABL.CSPInstance, FABL.UniformStringTester, FABL.UniformStringTester.constraint}
Let $D$ be a domain, let a finite predicate template specify each predicate’s arity and
satisfaction rule, and let a uniform string tester have a finite nonempty seed space.
The induced CSP instance contains, with multiplicity, one constraint for each seed,
using that seed’s predicate and ordered tuple of distinct variables; every declared
variable occurs in at least one such constraint.
\end{definition}

\begin{definition}[Tester acceptance probability]
\lean{FABL.UniformStringTester.acceptance}
\statementsource{Analysis-of-Boolean-Functions-by-Ryan-ODonnell}{pdf-ffbe9c7b3306-p187-b001}
\label{FABL.UniformStringTester.acceptance}
\leanok
\uses{FABL.UniformStringTester}
Let $D$ be a domain, let a finite predicate template specify a Boolean satisfaction rule
for each predicate, and let a uniform string tester have a finite nonempty seed space.
For an assignment of values in $D$ to every declared variable, its acceptance is the
probability, over a uniformly chosen seed, that the selected predicate is satisfied by
the assignment restricted to the ordered variables queried by that seed.
\end{definition}

\begin{theorem}[Tester acceptance equals instance value]
\lean{FABL.UniformStringTester.acceptance_eq_value}
\proofsource{Analysis-of-Boolean-Functions-by-Ryan-ODonnell}{pdf-ffbe9c7b3306-p187-b001}
\label{FABL.UniformStringTester.acceptance_eq_value}
\leanok
\uses{FABL.CSPConstraint.satisfiedIndicator, FABL.CSPInstance.satisfiedCount, FABL.CSPInstance.value, FABL.UniformStringTester, FABL.UniformStringTester.acceptance, FABL.UniformStringTester.constraint, FABL.UniformStringTester.toCSPInstance}
\difficulty{3}
\proofstep
  {FABL.𝔽₂}
  {direct}
  {Analysis-of-Boolean-Functions-by-Ryan-ODonnell}
  {pdf-ffbe9c7b3306-p015-b008}
\proofstep
  {FABL.Sign}
  {context}
  {Analysis-of-Boolean-Functions-by-Ryan-ODonnell}
  {pdf-ffbe9c7b3306-p014-b002,
    pdf-ffbe9c7b3306-p020-b001}
\proofstep
  {FABL.signValue}
  {context}
  {Analysis-of-Boolean-Functions-by-Ryan-ODonnell}
  {pdf-ffbe9c7b3306-p014-b002,
    pdf-ffbe9c7b3306-p018-b002}
\proofstep
  {FABL.signValue_one}
  {context}
  {Analysis-of-Boolean-Functions-by-Ryan-ODonnell}
  {pdf-ffbe9c7b3306-p018-b002}
\proofstep
  {FABL.signValue_neg_one}
  {context}
  {Analysis-of-Boolean-Functions-by-Ryan-ODonnell}
  {pdf-ffbe9c7b3306-p018-b002}
\proofstep
  {FABL.monomial}
  {direct}
  {Analysis-of-Boolean-Functions-by-Ryan-ODonnell}
  {pdf-ffbe9c7b3306-p022-b001,
    pdf-ffbe9c7b3306-p022-b004}
\proofstep
  {FABL.fourierCoeff}
  {direct}
  {Analysis-of-Boolean-Functions-by-Ryan-ODonnell}
  {pdf-ffbe9c7b3306-p025-b004}
\proofstep
  {FABL.fourierSupport}
  {context}
  {Analysis-of-Boolean-Functions-by-Ryan-ODonnell}
  {pdf-ffbe9c7b3306-p022-b002,
    pdf-ffbe9c7b3306-p072-b002}
\proofstep
  {FABL.mem_fourierSupport}
  {context}
  {Analysis-of-Boolean-Functions-by-Ryan-ODonnell}
  {pdf-ffbe9c7b3306-p072-b002}
\proofstep
  {FABL.granularityFrequencyWithin}
  {context}
  {Analysis-of-Boolean-Functions-by-Ryan-ODonnell}
  {pdf-ffbe9c7b3306-p076-b001}
\proofstep
  {FABL.mem_granularityFrequencyWithin}
  {context}
  {Analysis-of-Boolean-Functions-by-Ryan-ODonnell}
  {pdf-ffbe9c7b3306-p076-b001}
\proofstep
  {FABL.CSPTemplate}
  {context}
  {Analysis-of-Boolean-Functions-by-Ryan-ODonnell}
  {pdf-ffbe9c7b3306-p184-b004}
\proofstep
  {FABL.CSPTemplate.maxArity}
  {direct}
  {Analysis-of-Boolean-Functions-by-Ryan-ODonnell}
  {pdf-ffbe9c7b3306-p184-b004}
\proofstep
  {FABL.CSPTemplate.unsatisfiedCost}
  {context}
  {Analysis-of-Boolean-Functions-by-Ryan-ODonnell}
  {pdf-ffbe9c7b3306-p186-b001}
\proofstep
  {FABL.CSPTemplate.toValuedCSP}
  {context}
  {Analysis-of-Boolean-Functions-by-Ryan-ODonnell}
  {pdf-ffbe9c7b3306-p186-b001}
\proofstep
  {FABL.CSPTemplate.maxArity_eq_of_forall_arity_eq}
  {context}
  {Analysis-of-Boolean-Functions-by-Ryan-ODonnell}
  {pdf-ffbe9c7b3306-p184-b004}
\proofstep
  {FABL.CSPConstraint}
  {direct}
  {Analysis-of-Boolean-Functions-by-Ryan-ODonnell}
  {pdf-ffbe9c7b3306-p185-b004}
\proofstep
  {FABL.CSPConstraint.satisfiedIndicator}
  {context}
  {Analysis-of-Boolean-Functions-by-Ryan-ODonnell}
  {pdf-ffbe9c7b3306-p186-b001}
\proofstep
  {FABL.CSPConstraint.unsatisfiedIndicator}
  {context}
  {Analysis-of-Boolean-Functions-by-Ryan-ODonnell}
  {pdf-ffbe9c7b3306-p186-b001}
\proofstep
  {FABL.CSPConstraint.toValuedTerm}
  {context}
  {Analysis-of-Boolean-Functions-by-Ryan-ODonnell}
  {pdf-ffbe9c7b3306-p186-b001}
\proofstep
  {FABL.CSPConstraint.evalSolution_toValuedTerm}
  {context}
  {Analysis-of-Boolean-Functions-by-Ryan-ODonnell}
  {pdf-ffbe9c7b3306-p186-b001}
\proofstep
  {FABL.CSPInstance}
  {direct}
  {Analysis-of-Boolean-Functions-by-Ryan-ODonnell}
  {pdf-ffbe9c7b3306-p185-b004}
\proofstep
  {FABL.CSPInstance.satisfiedCount}
  {context}
  {Analysis-of-Boolean-Functions-by-Ryan-ODonnell}
  {pdf-ffbe9c7b3306-p186-b001}
\proofstep
  {FABL.CSPInstance.value}
  {direct}
  {Analysis-of-Boolean-Functions-by-Ryan-ODonnell}
  {pdf-ffbe9c7b3306-p185-b006,
    pdf-ffbe9c7b3306-p186-b001}
\proofstep
  {FABL.UniformStringTester}
  {context}
  {Analysis-of-Boolean-Functions-by-Ryan-ODonnell}
  {pdf-ffbe9c7b3306-p174-b001,
    pdf-ffbe9c7b3306-p184-b002}
\proofstep
  {FABL.UniformStringTester.constraint}
  {context}
  {Analysis-of-Boolean-Functions-by-Ryan-ODonnell}
  {pdf-ffbe9c7b3306-p185-b004}
\proofstep
  {FABL.UniformStringTester.toCSPInstance}
  {context}
  {Analysis-of-Boolean-Functions-by-Ryan-ODonnell}
  {pdf-ffbe9c7b3306-p187-b001}
\proofstep
  {FABL.UniformStringTester.acceptance}
  {direct}
  {Analysis-of-Boolean-Functions-by-Ryan-ODonnell}
  {pdf-ffbe9c7b3306-p174-b001}
\proofstep
  {FABL.maxCutTemplate}
  {context}
  {Analysis-of-Boolean-Functions-by-Ryan-ODonnell}
  {pdf-ffbe9c7b3306-p184-b003}
\proofstep
  {FABL.maxThreeColoringTemplate}
  {context}
  {Analysis-of-Boolean-Functions-by-Ryan-ODonnell}
  {pdf-ffbe9c7b3306-p184-b003}
\proofstep
  {FABL.maxE3LinTemplate}
  {context}
  {Analysis-of-Boolean-Functions-by-Ryan-ODonnell}
  {pdf-ffbe9c7b3306-p184-b003,
    pdf-ffbe9c7b3306-p185-b002}
\proofstep
  {FABL.e3SatPredicate}
  {context}
  {Analysis-of-Boolean-Functions-by-Ryan-ODonnell}
  {pdf-ffbe9c7b3306-p184-b003,
    pdf-ffbe9c7b3306-p185-b003}
\proofstep
  {FABL.maxE3SatTemplate}
  {context}
  {Analysis-of-Boolean-Functions-by-Ryan-ODonnell}
  {pdf-ffbe9c7b3306-p184-b003,
    pdf-ffbe9c7b3306-p185-b003}
\proofstep
  {FABL.maxArity_maxCutTemplate}
  {context}
  {Analysis-of-Boolean-Functions-by-Ryan-ODonnell}
  {pdf-ffbe9c7b3306-p184-b003,
    pdf-ffbe9c7b3306-p184-b004}
\proofstep
  {FABL.maxArity_maxThreeColoringTemplate}
  {context}
  {Analysis-of-Boolean-Functions-by-Ryan-ODonnell}
  {pdf-ffbe9c7b3306-p184-b003,
    pdf-ffbe9c7b3306-p184-b004}
\proofstep
  {FABL.maxArity_maxE3LinTemplate}
  {context}
  {Analysis-of-Boolean-Functions-by-Ryan-ODonnell}
  {pdf-ffbe9c7b3306-p184-b003,
    pdf-ffbe9c7b3306-p185-b002}
\proofstep
  {FABL.maxArity_maxE3SatTemplate}
  {context}
  {Analysis-of-Boolean-Functions-by-Ryan-ODonnell}
  {pdf-ffbe9c7b3306-p184-b003,
    pdf-ffbe9c7b3306-p185-b003}
Let $D$ be a domain, let $V$ be a set of variables, and let a finite family of
predicates on $D$ be specified, each with a finite arity. A tester chooses a predicate
and an ordered tuple of distinct variables of the corresponding arity using a uniform
draw from a finite nonempty seed set, and every variable in $V$ occurs in at least one
selected tuple. For every assignment of values in $D$ to $V$, the tester’s acceptance
probability equals the fraction of constraints satisfied in the associated instance,
where the instance counts each seed’s constraint with its multiplicity.
\end{theorem}

\begin{definition}[Max-cut predicate template]
\lean{FABL.maxCutTemplate}
\statementsource{Analysis-of-Boolean-Functions-by-Ryan-ODonnell}{pdf-ffbe9c7b3306-p184-b003}
\label{FABL.maxCutTemplate}
\leanok
\uses{FABL.CSPTemplate, FABL.Sign}
The max-cut template has a single binary predicate on signs in $bbZ^{\times}$, and
accepts exactly when its two inputs are unequal.
\end{definition}

\begin{definition}[Three-coloring predicate template]
\lean{FABL.maxThreeColoringTemplate}
\statementsource{Analysis-of-Boolean-Functions-by-Ryan-ODonnell}{pdf-ffbe9c7b3306-p184-b003}
\label{FABL.maxThreeColoringTemplate}
\leanok
\uses{FABL.CSPTemplate}
The max-three-coloring template has a single binary predicate on the three colors $bbN$
modulo $3$, and accepts exactly when its two inputs are unequal.
\end{definition}

\begin{definition}[E3-linear predicate template]
\lean{FABL.maxE3LinTemplate}
\statementsource{Analysis-of-Boolean-Functions-by-Ryan-ODonnell}{
  pdf-ffbe9c7b3306-p184-b003,
  pdf-ffbe9c7b3306-p185-b002,
  pdf-ffbe9c7b3306-p185-b003
}
\label{FABL.maxE3LinTemplate}
\leanok
\uses{FABL.CSPTemplate, FABL.𝔽₂}
The max-$3$-variable linear template over the field with two elements has one ternary
predicate for each right-hand side. It accepts a triple exactly when the sum of its
three entries equals that right-hand side.
\end{definition}

\begin{definition}[Three-literal satisfaction predicate]
\lean{FABL.e3SatPredicate}
\statementsource{Analysis-of-Boolean-Functions-by-Ryan-ODonnell}{
  pdf-ffbe9c7b3306-p184-b003,
  pdf-ffbe9c7b3306-p184-b005,
  pdf-ffbe9c7b3306-p185-b003
}
\label{FABL.e3SatPredicate}
\leanok
Given three Boolean negation flags and three Boolean values, the predicate is true
exactly when at least one of the three flags differs from its corresponding value.
\end{definition}

\begin{definition}[Coordinatewise negation equivalence]
\lean{FABL.e3SatNegationEquiv}
\label{FABL.e3SatNegationEquiv}
\leanok
For a fixed choice of which of three Boolean coordinates to negate, the map on Boolean
triples that flips exactly those chosen coordinates is a bijection from Boolean triples
to Boolean triples.
\end{definition}

\begin{theorem}[Three-literal satisfaction probability]
\lean{FABL.expect_plainE3SatIndicator}
\difficulty{2}
\proofsource{Analysis-of-Boolean-Functions-by-Ryan-ODonnell}{pdf-ffbe9c7b3306-p190-b004}
\proofstep
  {FABL.𝔽₂}
  {context}
  {Analysis-of-Boolean-Functions-by-Ryan-ODonnell}
  {pdf-ffbe9c7b3306-p015-b008,
    pdf-ffbe9c7b3306-p184-b003}
\proofstep
  {FABL.Sign}
  {context}
  {Analysis-of-Boolean-Functions-by-Ryan-ODonnell}
  {pdf-ffbe9c7b3306-p020-b001,
    pdf-ffbe9c7b3306-p022-b004}
\proofstep
  {FABL.signValue}
  {context}
  {Analysis-of-Boolean-Functions-by-Ryan-ODonnell}
  {pdf-ffbe9c7b3306-p022-b004}
\proofstep
  {FABL.signValue_one}
  {context}
  {Analysis-of-Boolean-Functions-by-Ryan-ODonnell}
  {pdf-ffbe9c7b3306-p022-b004}
\proofstep
  {FABL.signValue_neg_one}
  {context}
  {Analysis-of-Boolean-Functions-by-Ryan-ODonnell}
  {pdf-ffbe9c7b3306-p022-b004}
\proofstep
  {FABL.monomial}
  {direct}
  {Analysis-of-Boolean-Functions-by-Ryan-ODonnell}
  {pdf-ffbe9c7b3306-p022-b001,
    pdf-ffbe9c7b3306-p022-b004}
\proofstep
  {FABL.fourierCoeff}
  {direct}
  {Analysis-of-Boolean-Functions-by-Ryan-ODonnell}
  {pdf-ffbe9c7b3306-p025-b004}
\proofstep
  {FABL.fourierSupport}
  {context}
  {Analysis-of-Boolean-Functions-by-Ryan-ODonnell}
  {pdf-ffbe9c7b3306-p022-b002,
    pdf-ffbe9c7b3306-p018-b003}
\proofstep
  {FABL.mem_fourierSupport}
  {context}
  {Analysis-of-Boolean-Functions-by-Ryan-ODonnell}
  {pdf-ffbe9c7b3306-p018-b003}
\proofstep
  {FABL.granularityFrequencyWithin}
  {context}
  {Analysis-of-Boolean-Functions-by-Ryan-ODonnell}
  {pdf-ffbe9c7b3306-p072-b003}
\proofstep
  {FABL.mem_granularityFrequencyWithin}
  {context}
  {Analysis-of-Boolean-Functions-by-Ryan-ODonnell}
  {pdf-ffbe9c7b3306-p072-b003}
\proofstep
  {FABL.CSPTemplate}
  {direct}
  {Analysis-of-Boolean-Functions-by-Ryan-ODonnell}
  {pdf-ffbe9c7b3306-p184-b004}
\proofstep
  {FABL.CSPTemplate.maxArity}
  {direct}
  {Analysis-of-Boolean-Functions-by-Ryan-ODonnell}
  {pdf-ffbe9c7b3306-p184-b004}
\proofstep
  {FABL.CSPTemplate.unsatisfiedCost}
  {context}
  {Analysis-of-Boolean-Functions-by-Ryan-ODonnell}
  {pdf-ffbe9c7b3306-p186-b001}
\proofstep
  {FABL.CSPTemplate.toValuedCSP}
  {context}
  {Analysis-of-Boolean-Functions-by-Ryan-ODonnell}
  {pdf-ffbe9c7b3306-p184-b004,
    pdf-ffbe9c7b3306-p186-b001}
\proofstep
  {FABL.CSPTemplate.maxArity_eq_of_forall_arity_eq}
  {context}
  {Analysis-of-Boolean-Functions-by-Ryan-ODonnell}
  {pdf-ffbe9c7b3306-p184-b004}
\proofstep
  {FABL.CSPConstraint}
  {direct}
  {Analysis-of-Boolean-Functions-by-Ryan-ODonnell}
  {pdf-ffbe9c7b3306-p185-b004}
\proofstep
  {FABL.CSPConstraint.unsatisfiedIndicator}
  {context}
  {Analysis-of-Boolean-Functions-by-Ryan-ODonnell}
  {pdf-ffbe9c7b3306-p186-b001}
\proofstep
  {FABL.CSPConstraint.toValuedTerm}
  {context}
  {Analysis-of-Boolean-Functions-by-Ryan-ODonnell}
  {pdf-ffbe9c7b3306-p185-b004,
    pdf-ffbe9c7b3306-p186-b001}
\proofstep
  {FABL.CSPConstraint.evalSolution_toValuedTerm}
  {context}
  {Analysis-of-Boolean-Functions-by-Ryan-ODonnell}
  {pdf-ffbe9c7b3306-p186-b001}
\proofstep
  {FABL.maxCutTemplate}
  {direct}
  {Analysis-of-Boolean-Functions-by-Ryan-ODonnell}
  {pdf-ffbe9c7b3306-p184-b003}
\proofstep
  {FABL.maxThreeColoringTemplate}
  {direct}
  {Analysis-of-Boolean-Functions-by-Ryan-ODonnell}
  {pdf-ffbe9c7b3306-p184-b003}
\proofstep
  {FABL.maxE3LinTemplate}
  {direct}
  {Analysis-of-Boolean-Functions-by-Ryan-ODonnell}
  {pdf-ffbe9c7b3306-p184-b003,
    pdf-ffbe9c7b3306-p185-b002}
\proofstep
  {FABL.e3SatPredicate}
  {context}
  {Analysis-of-Boolean-Functions-by-Ryan-ODonnell}
  {pdf-ffbe9c7b3306-p184-b003,
    pdf-ffbe9c7b3306-p185-b003}
\proofstep
  {FABL.maxE3SatTemplate}
  {direct}
  {Analysis-of-Boolean-Functions-by-Ryan-ODonnell}
  {pdf-ffbe9c7b3306-p184-b003,
    pdf-ffbe9c7b3306-p185-b003}
\proofstep
  {FABL.maxArity_maxCutTemplate}
  {context}
  {Analysis-of-Boolean-Functions-by-Ryan-ODonnell}
  {pdf-ffbe9c7b3306-p184-b003}
\proofstep
  {FABL.maxArity_maxThreeColoringTemplate}
  {context}
  {Analysis-of-Boolean-Functions-by-Ryan-ODonnell}
  {pdf-ffbe9c7b3306-p184-b003}
\proofstep
  {FABL.maxArity_maxE3LinTemplate}
  {context}
  {Analysis-of-Boolean-Functions-by-Ryan-ODonnell}
  {pdf-ffbe9c7b3306-p184-b003,
    pdf-ffbe9c7b3306-p185-b002}
\proofstep
  {FABL.maxArity_maxE3SatTemplate}
  {context}
  {Analysis-of-Boolean-Functions-by-Ryan-ODonnell}
  {pdf-ffbe9c7b3306-p184-b003,
    pdf-ffbe9c7b3306-p185-b003}
\label{FABL.expect_plainE3SatIndicator}
\leanok
For three independent uniformly chosen Boolean values $v_0,v_1,v_2$, the probability
that at least one of them is true is $7/8$.
\end{theorem}

\begin{theorem}[Acceptance probability of a three-literal clause]
\lean{FABL.expect_e3SatPredicate}
\proofsource{Analysis-of-Boolean-Functions-by-Ryan-ODonnell}{
  pdf-ffbe9c7b3306-p190-b002,
  pdf-ffbe9c7b3306-p190-b004
}
\label{FABL.expect_e3SatPredicate}
\leanok
\uses{FABL.e3SatNegationEquiv, FABL.e3SatPredicate, FABL.expect_plainE3SatIndicator}
\difficulty{4}
\proofstep
  {FABL.𝔽₂}
  {context}
  {Analysis-of-Boolean-Functions-by-Ryan-ODonnell}
  {pdf-ffbe9c7b3306-p015-b008}
\proofstep
  {FABL.Sign}
  {context}
  {Analysis-of-Boolean-Functions-by-Ryan-ODonnell}
  {pdf-ffbe9c7b3306-p014-b002}
\proofstep
  {FABL.signValue}
  {context}
  {Analysis-of-Boolean-Functions-by-Ryan-ODonnell}
  {pdf-ffbe9c7b3306-p014-b002}
\proofstep
  {FABL.signValue_one}
  {context}
  {Analysis-of-Boolean-Functions-by-Ryan-ODonnell}
  {pdf-ffbe9c7b3306-p014-b002}
\proofstep
  {FABL.signValue_neg_one}
  {context}
  {Analysis-of-Boolean-Functions-by-Ryan-ODonnell}
  {pdf-ffbe9c7b3306-p014-b002}
\proofstep
  {FABL.monomial}
  {direct}
  {Analysis-of-Boolean-Functions-by-Ryan-ODonnell}
  {pdf-ffbe9c7b3306-p022-b001,
    pdf-ffbe9c7b3306-p022-b004}
\proofstep
  {FABL.fourierCoeff}
  {direct}
  {Analysis-of-Boolean-Functions-by-Ryan-ODonnell}
  {pdf-ffbe9c7b3306-p025-b004}
\proofstep
  {FABL.fourierSupport}
  {context}
  {Analysis-of-Boolean-Functions-by-Ryan-ODonnell}
  {pdf-ffbe9c7b3306-p022-b002}
\proofstep
  {FABL.mem_fourierSupport}
  {context}
  {Analysis-of-Boolean-Functions-by-Ryan-ODonnell}
  {pdf-ffbe9c7b3306-p022-b002}
\proofstep
  {FABL.granularityFrequencyWithin}
  {context}
  {Analysis-of-Boolean-Functions-by-Ryan-ODonnell}
  {pdf-ffbe9c7b3306-p072-b003}
\proofstep
  {FABL.mem_granularityFrequencyWithin}
  {context}
  {Analysis-of-Boolean-Functions-by-Ryan-ODonnell}
  {pdf-ffbe9c7b3306-p072-b003}
\proofstep
  {FABL.CSPTemplate}
  {context}
  {Analysis-of-Boolean-Functions-by-Ryan-ODonnell}
  {pdf-ffbe9c7b3306-p184-b004}
\proofstep
  {FABL.CSPTemplate.maxArity}
  {direct}
  {Analysis-of-Boolean-Functions-by-Ryan-ODonnell}
  {pdf-ffbe9c7b3306-p184-b004}
\proofstep
  {FABL.CSPTemplate.unsatisfiedCost}
  {context}
  {Analysis-of-Boolean-Functions-by-Ryan-ODonnell}
  {pdf-ffbe9c7b3306-p186-b001}
\proofstep
  {FABL.CSPTemplate.toValuedCSP}
  {context}
  {Analysis-of-Boolean-Functions-by-Ryan-ODonnell}
  {pdf-ffbe9c7b3306-p184-b004,
    pdf-ffbe9c7b3306-p186-b001}
\proofstep
  {FABL.CSPTemplate.maxArity_eq_of_forall_arity_eq}
  {context}
  {Analysis-of-Boolean-Functions-by-Ryan-ODonnell}
  {pdf-ffbe9c7b3306-p184-b004}
\proofstep
  {FABL.CSPConstraint}
  {direct}
  {Analysis-of-Boolean-Functions-by-Ryan-ODonnell}
  {pdf-ffbe9c7b3306-p185-b004}
\proofstep
  {FABL.CSPConstraint.unsatisfiedIndicator}
  {context}
  {Analysis-of-Boolean-Functions-by-Ryan-ODonnell}
  {pdf-ffbe9c7b3306-p186-b001}
\proofstep
  {FABL.CSPConstraint.toValuedTerm}
  {context}
  {Analysis-of-Boolean-Functions-by-Ryan-ODonnell}
  {pdf-ffbe9c7b3306-p185-b004,
    pdf-ffbe9c7b3306-p186-b001}
\proofstep
  {FABL.CSPConstraint.evalSolution_toValuedTerm}
  {context}
  {Analysis-of-Boolean-Functions-by-Ryan-ODonnell}
  {pdf-ffbe9c7b3306-p186-b001}
\proofstep
  {FABL.maxCutTemplate}
  {context}
  {Analysis-of-Boolean-Functions-by-Ryan-ODonnell}
  {pdf-ffbe9c7b3306-p184-b003,
    pdf-ffbe9c7b3306-p184-b004}
\proofstep
  {FABL.maxThreeColoringTemplate}
  {context}
  {Analysis-of-Boolean-Functions-by-Ryan-ODonnell}
  {pdf-ffbe9c7b3306-p184-b003,
    pdf-ffbe9c7b3306-p184-b004}
\proofstep
  {FABL.maxE3LinTemplate}
  {context}
  {Analysis-of-Boolean-Functions-by-Ryan-ODonnell}
  {pdf-ffbe9c7b3306-p184-b003,
    pdf-ffbe9c7b3306-p185-b002}
\proofstep
  {FABL.e3SatPredicate}
  {context}
  {Analysis-of-Boolean-Functions-by-Ryan-ODonnell}
  {pdf-ffbe9c7b3306-p190-b002,
    pdf-ffbe9c7b3306-p190-b004}
\proofstep
  {FABL.e3SatNegationEquiv}
  {context}
  {Analysis-of-Boolean-Functions-by-Ryan-ODonnell}
  {pdf-ffbe9c7b3306-p190-b002}
\proofstep
  {FABL.expect_plainE3SatIndicator}
  {direct}
  {Analysis-of-Boolean-Functions-by-Ryan-ODonnell}
  {pdf-ffbe9c7b3306-p190-b002,
    pdf-ffbe9c7b3306-p190-b004}
\proofstep
  {FABL.maxE3SatTemplate}
  {context}
  {Analysis-of-Boolean-Functions-by-Ryan-ODonnell}
  {pdf-ffbe9c7b3306-p184-b003,
    pdf-ffbe9c7b3306-p190-b002}
\proofstep
  {FABL.maxArity_maxCutTemplate}
  {context}
  {Analysis-of-Boolean-Functions-by-Ryan-ODonnell}
  {pdf-ffbe9c7b3306-p184-b003,
    pdf-ffbe9c7b3306-p184-b004}
\proofstep
  {FABL.maxArity_maxThreeColoringTemplate}
  {context}
  {Analysis-of-Boolean-Functions-by-Ryan-ODonnell}
  {pdf-ffbe9c7b3306-p184-b003,
    pdf-ffbe9c7b3306-p184-b004}
\proofstep
  {FABL.maxArity_maxE3LinTemplate}
  {context}
  {Analysis-of-Boolean-Functions-by-Ryan-ODonnell}
  {pdf-ffbe9c7b3306-p184-b004,
    pdf-ffbe9c7b3306-p185-b002}
\proofstep
  {FABL.maxArity_maxE3SatTemplate}
  {context}
  {Analysis-of-Boolean-Functions-by-Ryan-ODonnell}
  {pdf-ffbe9c7b3306-p184-b004,
    pdf-ffbe9c7b3306-p190-b002}
Fix any three Boolean negation choices. A uniformly random Boolean assignment to the
three variables satisfies the resulting disjunction of literals with probability $7/8$.
\end{theorem}

\begin{definition}[Max-three-SAT predicate template]
\lean{FABL.maxE3SatTemplate}
\statementsource{Analysis-of-Boolean-Functions-by-Ryan-ODonnell}{
  pdf-ffbe9c7b3306-p184-b003,
  pdf-ffbe9c7b3306-p184-b005,
  pdf-ffbe9c7b3306-p185-b003
}
\label{FABL.maxE3SatTemplate}
\leanok
\uses{FABL.CSPTemplate, FABL.e3SatPredicate}
The max-three-SAT template has a predicate for each triple of Boolean negation flags and
has arity three. On a triple of Boolean inputs, it evaluates the three-literal
satisfaction predicate: it is true exactly when at least one negation flag differs from
its corresponding input.
\end{definition}

\begin{definition}[Assignment split by a constraint scope]
\lean{FABL.e3SatScopeSplitEquiv}
\label{FABL.e3SatScopeSplitEquiv}
\leanok
\uses{FABL.CSPConstraint, FABL.maxE3SatTemplate}
Let $V$ be a finite variable set and let $c$ be a constraint of arity three on $V$. An
assignment of Boolean values to $V$ is equivalent to a pair consisting of its values on
the three variables in the ordered scope of $c$ and its values on all variables outside
that scope.
\end{definition}

\begin{theorem}[Constraint assignment on its scope]
\lean{FABL.e3SatScopeSplitEquiv_fst_apply}
\label{FABL.e3SatScopeSplitEquiv_fst_apply}
\leanok
\uses{FABL.CSPConstraint, FABL.e3SatScopeSplitEquiv, FABL.maxE3SatTemplate}
\difficulty{3}
Let $V$ be a finite set of variables, and let $c$ be a three-literal satisfiability
constraint whose scope is an ordered triple of distinct variables. For any Boolean
assignment $a:V\to\{\mathrm{false},\mathrm{true}\}$, the first component of the
assignment split into its values on the scope and its values off the scope is the
function on the three scope positions given by $i\mapsto a(c_i)$.
\end{theorem}

\begin{theorem}[Expected satisfaction of a three-literal clause]
\lean{FABL.expect_satisfiedIndicator_maxE3Sat}
\proofsource{Analysis-of-Boolean-Functions-by-Ryan-ODonnell}{
  pdf-ffbe9c7b3306-p190-b002,
  pdf-ffbe9c7b3306-p190-b004
}
\label{FABL.expect_satisfiedIndicator_maxE3Sat}
\leanok
\uses{FABL.CSPConstraint, FABL.CSPConstraint.satisfiedIndicator, FABL.e3SatPredicate, FABL.e3SatScopeSplitEquiv, FABL.e3SatScopeSplitEquiv_fst_apply, FABL.expect_e3SatPredicate, FABL.maxE3SatTemplate}
\difficulty{5}
\proofstep
  {FABL.𝔽₂}
  {direct}
  {Analysis-of-Boolean-Functions-by-Ryan-ODonnell}
  {pdf-ffbe9c7b3306-p015-b008}
\proofstep
  {FABL.Sign}
  {context}
  {Analysis-of-Boolean-Functions-by-Ryan-ODonnell}
  {pdf-ffbe9c7b3306-p020-b001,
    pdf-ffbe9c7b3306-p014-b002}
\proofstep
  {FABL.signValue}
  {context}
  {Analysis-of-Boolean-Functions-by-Ryan-ODonnell}
  {pdf-ffbe9c7b3306-p014-b002}
\proofstep
  {FABL.signValue_one}
  {context}
  {Analysis-of-Boolean-Functions-by-Ryan-ODonnell}
  {pdf-ffbe9c7b3306-p014-b002}
\proofstep
  {FABL.signValue_neg_one}
  {context}
  {Analysis-of-Boolean-Functions-by-Ryan-ODonnell}
  {pdf-ffbe9c7b3306-p014-b002}
\proofstep
  {FABL.monomial}
  {direct}
  {Analysis-of-Boolean-Functions-by-Ryan-ODonnell}
  {pdf-ffbe9c7b3306-p022-b001,
    pdf-ffbe9c7b3306-p022-b004}
\proofstep
  {FABL.fourierCoeff}
  {direct}
  {Analysis-of-Boolean-Functions-by-Ryan-ODonnell}
  {pdf-ffbe9c7b3306-p025-b004}
\proofstep
  {FABL.fourierSupport}
  {context}
  {Analysis-of-Boolean-Functions-by-Ryan-ODonnell}
  {pdf-ffbe9c7b3306-p022-b002,
    pdf-ffbe9c7b3306-p022-b001}
\proofstep
  {FABL.mem_fourierSupport}
  {context}
  {Analysis-of-Boolean-Functions-by-Ryan-ODonnell}
  {pdf-ffbe9c7b3306-p022-b002}
\proofstep
  {FABL.granularityFrequencyWithin}
  {context}
  {Analysis-of-Boolean-Functions-by-Ryan-ODonnell}
  {pdf-ffbe9c7b3306-p015-b012}
\proofstep
  {FABL.mem_granularityFrequencyWithin}
  {context}
  {Analysis-of-Boolean-Functions-by-Ryan-ODonnell}
  {pdf-ffbe9c7b3306-p015-b012}
\proofstep
  {FABL.CSPTemplate}
  {direct}
  {Analysis-of-Boolean-Functions-by-Ryan-ODonnell}
  {pdf-ffbe9c7b3306-p184-b004}
\proofstep
  {FABL.CSPTemplate.maxArity}
  {direct}
  {Analysis-of-Boolean-Functions-by-Ryan-ODonnell}
  {pdf-ffbe9c7b3306-p184-b004}
\proofstep
  {FABL.CSPTemplate.unsatisfiedCost}
  {context}
  {Analysis-of-Boolean-Functions-by-Ryan-ODonnell}
  {pdf-ffbe9c7b3306-p186-b001}
\proofstep
  {FABL.CSPTemplate.toValuedCSP}
  {context}
  {Analysis-of-Boolean-Functions-by-Ryan-ODonnell}
  {pdf-ffbe9c7b3306-p186-b001,
    pdf-ffbe9c7b3306-p185-b006}
\proofstep
  {FABL.CSPTemplate.maxArity_eq_of_forall_arity_eq}
  {context}
  {Analysis-of-Boolean-Functions-by-Ryan-ODonnell}
  {pdf-ffbe9c7b3306-p184-b004}
\proofstep
  {FABL.CSPConstraint}
  {direct}
  {Analysis-of-Boolean-Functions-by-Ryan-ODonnell}
  {pdf-ffbe9c7b3306-p185-b004}
\proofstep
  {FABL.CSPConstraint.satisfiedIndicator}
  {context}
  {Analysis-of-Boolean-Functions-by-Ryan-ODonnell}
  {pdf-ffbe9c7b3306-p186-b001}
\proofstep
  {FABL.CSPConstraint.unsatisfiedIndicator}
  {context}
  {Analysis-of-Boolean-Functions-by-Ryan-ODonnell}
  {pdf-ffbe9c7b3306-p186-b001}
\proofstep
  {FABL.CSPConstraint.toValuedTerm}
  {context}
  {Analysis-of-Boolean-Functions-by-Ryan-ODonnell}
  {pdf-ffbe9c7b3306-p186-b001,
    pdf-ffbe9c7b3306-p185-b004}
\proofstep
  {FABL.CSPConstraint.evalSolution_toValuedTerm}
  {context}
  {Analysis-of-Boolean-Functions-by-Ryan-ODonnell}
  {pdf-ffbe9c7b3306-p186-b001}
\proofstep
  {FABL.maxCutTemplate}
  {direct}
  {Analysis-of-Boolean-Functions-by-Ryan-ODonnell}
  {pdf-ffbe9c7b3306-p184-b003}
\proofstep
  {FABL.maxThreeColoringTemplate}
  {direct}
  {Analysis-of-Boolean-Functions-by-Ryan-ODonnell}
  {pdf-ffbe9c7b3306-p184-b003}
\proofstep
  {FABL.maxE3LinTemplate}
  {direct}
  {Analysis-of-Boolean-Functions-by-Ryan-ODonnell}
  {pdf-ffbe9c7b3306-p184-b003,
    pdf-ffbe9c7b3306-p185-b002}
\proofstep
  {FABL.e3SatPredicate}
  {context}
  {Analysis-of-Boolean-Functions-by-Ryan-ODonnell}
  {pdf-ffbe9c7b3306-p190-b002,
    pdf-ffbe9c7b3306-p190-b004}
\proofstep
  {FABL.e3SatNegationEquiv}
  {context}
  {Analysis-of-Boolean-Functions-by-Ryan-ODonnell}
  {pdf-ffbe9c7b3306-p190-b004}
\proofstep
  {FABL.expect_plainE3SatIndicator}
  {direct}
  {Analysis-of-Boolean-Functions-by-Ryan-ODonnell}
  {pdf-ffbe9c7b3306-p190-b002,
    pdf-ffbe9c7b3306-p190-b004}
\proofstep
  {FABL.expect_e3SatPredicate}
  {direct}
  {Analysis-of-Boolean-Functions-by-Ryan-ODonnell}
  {pdf-ffbe9c7b3306-p190-b004}
\proofstep
  {FABL.maxE3SatTemplate}
  {direct}
  {Analysis-of-Boolean-Functions-by-Ryan-ODonnell}
  {pdf-ffbe9c7b3306-p184-b003,
    pdf-ffbe9c7b3306-p190-b004}
\proofstep
  {FABL.e3SatScopeSplitEquiv}
  {context}
  {Analysis-of-Boolean-Functions-by-Ryan-ODonnell}
  {pdf-ffbe9c7b3306-p185-b004}
\proofstep
  {FABL.e3SatScopeSplitEquiv_fst_apply}
  {context}
  {Analysis-of-Boolean-Functions-by-Ryan-ODonnell}
  {pdf-ffbe9c7b3306-p185-b004}
\proofstep
  {FABL.maxArity_maxCutTemplate}
  {context}
  {Analysis-of-Boolean-Functions-by-Ryan-ODonnell}
  {pdf-ffbe9c7b3306-p184-b003,
    pdf-ffbe9c7b3306-p184-b004}
\proofstep
  {FABL.maxArity_maxThreeColoringTemplate}
  {context}
  {Analysis-of-Boolean-Functions-by-Ryan-ODonnell}
  {pdf-ffbe9c7b3306-p184-b003,
    pdf-ffbe9c7b3306-p184-b004}
\proofstep
  {FABL.maxArity_maxE3LinTemplate}
  {context}
  {Analysis-of-Boolean-Functions-by-Ryan-ODonnell}
  {pdf-ffbe9c7b3306-p184-b003,
    pdf-ffbe9c7b3306-p185-b002}
\proofstep
  {FABL.maxArity_maxE3SatTemplate}
  {context}
  {Analysis-of-Boolean-Functions-by-Ryan-ODonnell}
  {pdf-ffbe9c7b3306-p184-b003,
    pdf-ffbe9c7b3306-p190-b004}
Let $V$ be a finite set and consider a clause on three distinct variables, each
appearing either positively or negated, which is satisfied when at least one of its
three literals is true. Under a uniformly random Boolean assignment to $V$, the expected
satisfaction indicator of the clause is $7/8$.
\end{theorem}

\begin{theorem}[Expected number of satisfied three-literal clauses]
\lean{FABL.expect_satisfiedIndicators_maxE3Sat}
\proofsource{Analysis-of-Boolean-Functions-by-Ryan-ODonnell}{
  pdf-ffbe9c7b3306-p190-b002,
  pdf-ffbe9c7b3306-p190-b004
}
\label{FABL.expect_satisfiedIndicators_maxE3Sat}
\leanok
\uses{FABL.CSPConstraint, FABL.CSPConstraint.satisfiedIndicator, FABL.expect_satisfiedIndicator_maxE3Sat, FABL.maxE3SatTemplate}
\difficulty{4}
\proofstep
  {FABL.𝔽₂}
  {direct}
  {Analysis-of-Boolean-Functions-by-Ryan-ODonnell}
  {pdf-ffbe9c7b3306-p015-b008}
\proofstep
  {FABL.Sign}
  {context}
  {Analysis-of-Boolean-Functions-by-Ryan-ODonnell}
  {pdf-ffbe9c7b3306-p020-b001,
    pdf-ffbe9c7b3306-p014-b002}
\proofstep
  {FABL.signValue}
  {context}
  {Analysis-of-Boolean-Functions-by-Ryan-ODonnell}
  {pdf-ffbe9c7b3306-p014-b002}
\proofstep
  {FABL.signValue_one}
  {context}
  {Analysis-of-Boolean-Functions-by-Ryan-ODonnell}
  {pdf-ffbe9c7b3306-p014-b002}
\proofstep
  {FABL.signValue_neg_one}
  {context}
  {Analysis-of-Boolean-Functions-by-Ryan-ODonnell}
  {pdf-ffbe9c7b3306-p014-b002}
\proofstep
  {FABL.monomial}
  {direct}
  {Analysis-of-Boolean-Functions-by-Ryan-ODonnell}
  {pdf-ffbe9c7b3306-p014-b002,
    pdf-ffbe9c7b3306-p022-b001}
\proofstep
  {FABL.fourierCoeff}
  {direct}
  {Analysis-of-Boolean-Functions-by-Ryan-ODonnell}
  {pdf-ffbe9c7b3306-p025-b004}
\proofstep
  {FABL.fourierSupport}
  {context}
  {Analysis-of-Boolean-Functions-by-Ryan-ODonnell}
  {pdf-ffbe9c7b3306-p022-b002}
\proofstep
  {FABL.mem_fourierSupport}
  {context}
  {Analysis-of-Boolean-Functions-by-Ryan-ODonnell}
  {pdf-ffbe9c7b3306-p022-b002}
\proofstep
  {FABL.granularityFrequencyWithin}
  {context}
  {Analysis-of-Boolean-Functions-by-Ryan-ODonnell}
  {pdf-ffbe9c7b3306-p013-b007}
\proofstep
  {FABL.mem_granularityFrequencyWithin}
  {context}
  {Analysis-of-Boolean-Functions-by-Ryan-ODonnell}
  {pdf-ffbe9c7b3306-p013-b007}
\proofstep
  {FABL.CSPTemplate}
  {direct}
  {Analysis-of-Boolean-Functions-by-Ryan-ODonnell}
  {pdf-ffbe9c7b3306-p184-b004}
\proofstep
  {FABL.CSPTemplate.maxArity}
  {direct}
  {Analysis-of-Boolean-Functions-by-Ryan-ODonnell}
  {pdf-ffbe9c7b3306-p184-b004}
\proofstep
  {FABL.CSPTemplate.unsatisfiedCost}
  {context}
  {Analysis-of-Boolean-Functions-by-Ryan-ODonnell}
  {pdf-ffbe9c7b3306-p186-b001}
\proofstep
  {FABL.CSPTemplate.toValuedCSP}
  {context}
  {Analysis-of-Boolean-Functions-by-Ryan-ODonnell}
  {pdf-ffbe9c7b3306-p184-b004,
    pdf-ffbe9c7b3306-p186-b001}
\proofstep
  {FABL.CSPTemplate.maxArity_eq_of_forall_arity_eq}
  {context}
  {Analysis-of-Boolean-Functions-by-Ryan-ODonnell}
  {pdf-ffbe9c7b3306-p184-b004}
\proofstep
  {FABL.CSPConstraint}
  {direct}
  {Analysis-of-Boolean-Functions-by-Ryan-ODonnell}
  {pdf-ffbe9c7b3306-p185-b004}
\proofstep
  {FABL.CSPConstraint.satisfiedIndicator}
  {direct}
  {Analysis-of-Boolean-Functions-by-Ryan-ODonnell}
  {pdf-ffbe9c7b3306-p186-b001}
\proofstep
  {FABL.CSPConstraint.unsatisfiedIndicator}
  {context}
  {Analysis-of-Boolean-Functions-by-Ryan-ODonnell}
  {pdf-ffbe9c7b3306-p186-b001}
\proofstep
  {FABL.CSPConstraint.toValuedTerm}
  {context}
  {Analysis-of-Boolean-Functions-by-Ryan-ODonnell}
  {pdf-ffbe9c7b3306-p185-b004,
    pdf-ffbe9c7b3306-p186-b001}
\proofstep
  {FABL.CSPConstraint.evalSolution_toValuedTerm}
  {context}
  {Analysis-of-Boolean-Functions-by-Ryan-ODonnell}
  {pdf-ffbe9c7b3306-p186-b001}
\proofstep
  {FABL.maxCutTemplate}
  {context}
  {Analysis-of-Boolean-Functions-by-Ryan-ODonnell}
  {pdf-ffbe9c7b3306-p184-b003}
\proofstep
  {FABL.maxThreeColoringTemplate}
  {context}
  {Analysis-of-Boolean-Functions-by-Ryan-ODonnell}
  {pdf-ffbe9c7b3306-p184-b003}
\proofstep
  {FABL.maxE3LinTemplate}
  {context}
  {Analysis-of-Boolean-Functions-by-Ryan-ODonnell}
  {pdf-ffbe9c7b3306-p184-b003,
    pdf-ffbe9c7b3306-p185-b002}
\proofstep
  {FABL.e3SatPredicate}
  {context}
  {Analysis-of-Boolean-Functions-by-Ryan-ODonnell}
  {pdf-ffbe9c7b3306-p184-b003}
\proofstep
  {FABL.e3SatNegationEquiv}
  {context}
  {Analysis-of-Boolean-Functions-by-Ryan-ODonnell}
  {pdf-ffbe9c7b3306-p184-b003}
\proofstep
  {FABL.expect_plainE3SatIndicator}
  {direct}
  {Analysis-of-Boolean-Functions-by-Ryan-ODonnell}
  {pdf-ffbe9c7b3306-p190-b002,
    pdf-ffbe9c7b3306-p190-b004}
\proofstep
  {FABL.expect_e3SatPredicate}
  {direct}
  {Analysis-of-Boolean-Functions-by-Ryan-ODonnell}
  {pdf-ffbe9c7b3306-p190-b002,
    pdf-ffbe9c7b3306-p190-b004}
\proofstep
  {FABL.maxE3SatTemplate}
  {context}
  {Analysis-of-Boolean-Functions-by-Ryan-ODonnell}
  {pdf-ffbe9c7b3306-p184-b003}
\proofstep
  {FABL.e3SatScopeSplitEquiv}
  {context}
  {Analysis-of-Boolean-Functions-by-Ryan-ODonnell}
  {pdf-ffbe9c7b3306-p186-b001}
\proofstep
  {FABL.e3SatScopeSplitEquiv_fst_apply}
  {context}
  {Analysis-of-Boolean-Functions-by-Ryan-ODonnell}
  {pdf-ffbe9c7b3306-p186-b001}
\proofstep
  {FABL.expect_satisfiedIndicator_maxE3Sat}
  {direct}
  {Analysis-of-Boolean-Functions-by-Ryan-ODonnell}
  {pdf-ffbe9c7b3306-p190-b002,
    pdf-ffbe9c7b3306-p190-b004}
\proofstep
  {FABL.maxArity_maxCutTemplate}
  {context}
  {Analysis-of-Boolean-Functions-by-Ryan-ODonnell}
  {pdf-ffbe9c7b3306-p184-b003,
    pdf-ffbe9c7b3306-p184-b004}
\proofstep
  {FABL.maxArity_maxThreeColoringTemplate}
  {context}
  {Analysis-of-Boolean-Functions-by-Ryan-ODonnell}
  {pdf-ffbe9c7b3306-p184-b003,
    pdf-ffbe9c7b3306-p184-b004}
\proofstep
  {FABL.maxArity_maxE3LinTemplate}
  {context}
  {Analysis-of-Boolean-Functions-by-Ryan-ODonnell}
  {pdf-ffbe9c7b3306-p184-b003,
    pdf-ffbe9c7b3306-p185-b002,
    pdf-ffbe9c7b3306-p184-b004}
\proofstep
  {FABL.maxArity_maxE3SatTemplate}
  {context}
  {Analysis-of-Boolean-Functions-by-Ryan-ODonnell}
  {pdf-ffbe9c7b3306-p184-b003,
    pdf-ffbe9c7b3306-p184-b004}
Let $V$ be a finite set and let $\mathcal{C}$ be a finite multiset of clauses, each on
three distinct variables, with each variable appearing either positively or negated. A
clause is satisfied when at least one of its literals is true. Under a uniformly random
Boolean assignment to $V$, the expected number of satisfied clauses in $\mathcal{C}$ is
$7|\mathcal{C}|/8$.
\end{theorem}

\begin{theorem}[Expected value of a three-literal clause instance]
\lean{FABL.expect_value_maxE3Sat}
\proofsource{Analysis-of-Boolean-Functions-by-Ryan-ODonnell}{
  pdf-ffbe9c7b3306-p190-b002,
  pdf-ffbe9c7b3306-p190-b004
}
\label{FABL.expect_value_maxE3Sat}
\leanok
\uses{FABL.CSPInstance, FABL.CSPInstance.satisfiedCount, FABL.CSPInstance.value, FABL.expect_satisfiedIndicators_maxE3Sat, FABL.maxE3SatTemplate}
\difficulty{4}
\proofstep
  {FABL.𝔽₂}
  {direct}
  {Analysis-of-Boolean-Functions-by-Ryan-ODonnell}
  {pdf-ffbe9c7b3306-p015-b008}
\proofstep
  {FABL.Sign}
  {context}
  {Analysis-of-Boolean-Functions-by-Ryan-ODonnell}
  {pdf-ffbe9c7b3306-p014-b002}
\proofstep
  {FABL.signValue}
  {context}
  {Analysis-of-Boolean-Functions-by-Ryan-ODonnell}
  {pdf-ffbe9c7b3306-p014-b002}
\proofstep
  {FABL.signValue_one}
  {context}
  {Analysis-of-Boolean-Functions-by-Ryan-ODonnell}
  {pdf-ffbe9c7b3306-p014-b002}
\proofstep
  {FABL.signValue_neg_one}
  {context}
  {Analysis-of-Boolean-Functions-by-Ryan-ODonnell}
  {pdf-ffbe9c7b3306-p014-b002}
\proofstep
  {FABL.monomial}
  {direct}
  {Analysis-of-Boolean-Functions-by-Ryan-ODonnell}
  {pdf-ffbe9c7b3306-p022-b001,
    pdf-ffbe9c7b3306-p022-b004}
\proofstep
  {FABL.fourierCoeff}
  {direct}
  {Analysis-of-Boolean-Functions-by-Ryan-ODonnell}
  {pdf-ffbe9c7b3306-p025-b004}
\proofstep
  {FABL.fourierSupport}
  {context}
  {Analysis-of-Boolean-Functions-by-Ryan-ODonnell}
  {pdf-ffbe9c7b3306-p018-b003,
    pdf-ffbe9c7b3306-p022-b002}
\proofstep
  {FABL.mem_fourierSupport}
  {context}
  {Analysis-of-Boolean-Functions-by-Ryan-ODonnell}
  {pdf-ffbe9c7b3306-p018-b003,
    pdf-ffbe9c7b3306-p022-b002}
\proofstep
  {FABL.granularityFrequencyWithin}
  {context}
  {Analysis-of-Boolean-Functions-by-Ryan-ODonnell}
  {pdf-ffbe9c7b3306-p076-b001,
    pdf-ffbe9c7b3306-p084-b002}
\proofstep
  {FABL.mem_granularityFrequencyWithin}
  {context}
  {Analysis-of-Boolean-Functions-by-Ryan-ODonnell}
  {pdf-ffbe9c7b3306-p076-b001,
    pdf-ffbe9c7b3306-p084-b002}
\proofstep
  {FABL.CSPTemplate}
  {direct}
  {Analysis-of-Boolean-Functions-by-Ryan-ODonnell}
  {pdf-ffbe9c7b3306-p184-b004}
\proofstep
  {FABL.CSPTemplate.maxArity}
  {direct}
  {Analysis-of-Boolean-Functions-by-Ryan-ODonnell}
  {pdf-ffbe9c7b3306-p184-b004}
\proofstep
  {FABL.CSPTemplate.unsatisfiedCost}
  {context}
  {Analysis-of-Boolean-Functions-by-Ryan-ODonnell}
  {pdf-ffbe9c7b3306-p186-b001}
\proofstep
  {FABL.CSPTemplate.toValuedCSP}
  {context}
  {Analysis-of-Boolean-Functions-by-Ryan-ODonnell}
  {pdf-ffbe9c7b3306-p186-b001}
\proofstep
  {FABL.CSPTemplate.maxArity_eq_of_forall_arity_eq}
  {context}
  {Analysis-of-Boolean-Functions-by-Ryan-ODonnell}
  {pdf-ffbe9c7b3306-p184-b004}
\proofstep
  {FABL.CSPConstraint}
  {direct}
  {Analysis-of-Boolean-Functions-by-Ryan-ODonnell}
  {pdf-ffbe9c7b3306-p185-b004}
\proofstep
  {FABL.CSPConstraint.satisfiedIndicator}
  {direct}
  {Analysis-of-Boolean-Functions-by-Ryan-ODonnell}
  {pdf-ffbe9c7b3306-p186-b001}
\proofstep
  {FABL.CSPConstraint.unsatisfiedIndicator}
  {context}
  {Analysis-of-Boolean-Functions-by-Ryan-ODonnell}
  {pdf-ffbe9c7b3306-p186-b001}
\proofstep
  {FABL.CSPConstraint.toValuedTerm}
  {context}
  {Analysis-of-Boolean-Functions-by-Ryan-ODonnell}
  {pdf-ffbe9c7b3306-p185-b004,
    pdf-ffbe9c7b3306-p186-b001}
\proofstep
  {FABL.CSPConstraint.evalSolution_toValuedTerm}
  {context}
  {Analysis-of-Boolean-Functions-by-Ryan-ODonnell}
  {pdf-ffbe9c7b3306-p186-b001}
\proofstep
  {FABL.CSPInstance}
  {direct}
  {Analysis-of-Boolean-Functions-by-Ryan-ODonnell}
  {pdf-ffbe9c7b3306-p185-b004}
\proofstep
  {FABL.CSPInstance.satisfiedCount}
  {context}
  {Analysis-of-Boolean-Functions-by-Ryan-ODonnell}
  {pdf-ffbe9c7b3306-p186-b001}
\proofstep
  {FABL.CSPInstance.value}
  {direct}
  {Analysis-of-Boolean-Functions-by-Ryan-ODonnell}
  {pdf-ffbe9c7b3306-p185-b006,
    pdf-ffbe9c7b3306-p186-b001}
\proofstep
  {FABL.maxCutTemplate}
  {direct}
  {Analysis-of-Boolean-Functions-by-Ryan-ODonnell}
  {pdf-ffbe9c7b3306-p184-b003}
\proofstep
  {FABL.maxThreeColoringTemplate}
  {direct}
  {Analysis-of-Boolean-Functions-by-Ryan-ODonnell}
  {pdf-ffbe9c7b3306-p184-b003}
\proofstep
  {FABL.maxE3LinTemplate}
  {direct}
  {Analysis-of-Boolean-Functions-by-Ryan-ODonnell}
  {pdf-ffbe9c7b3306-p184-b003,
    pdf-ffbe9c7b3306-p185-b002}
\proofstep
  {FABL.e3SatPredicate}
  {context}
  {Analysis-of-Boolean-Functions-by-Ryan-ODonnell}
  {pdf-ffbe9c7b3306-p184-b003,
    pdf-ffbe9c7b3306-p190-b004}
\proofstep
  {FABL.e3SatNegationEquiv}
  {context}
  {Analysis-of-Boolean-Functions-by-Ryan-ODonnell}
  {pdf-ffbe9c7b3306-p190-b002,
    pdf-ffbe9c7b3306-p190-b004}
\proofstep
  {FABL.expect_plainE3SatIndicator}
  {direct}
  {Analysis-of-Boolean-Functions-by-Ryan-ODonnell}
  {pdf-ffbe9c7b3306-p190-b002,
    pdf-ffbe9c7b3306-p190-b004}
\proofstep
  {FABL.expect_e3SatPredicate}
  {direct}
  {Analysis-of-Boolean-Functions-by-Ryan-ODonnell}
  {pdf-ffbe9c7b3306-p190-b002,
    pdf-ffbe9c7b3306-p190-b004}
\proofstep
  {FABL.maxE3SatTemplate}
  {direct}
  {Analysis-of-Boolean-Functions-by-Ryan-ODonnell}
  {pdf-ffbe9c7b3306-p184-b003,
    pdf-ffbe9c7b3306-p190-b002}
\proofstep
  {FABL.e3SatScopeSplitEquiv}
  {context}
  {Analysis-of-Boolean-Functions-by-Ryan-ODonnell}
  {pdf-ffbe9c7b3306-p185-b004,
    pdf-ffbe9c7b3306-p186-b001}
\proofstep
  {FABL.e3SatScopeSplitEquiv_fst_apply}
  {context}
  {Analysis-of-Boolean-Functions-by-Ryan-ODonnell}
  {pdf-ffbe9c7b3306-p185-b004,
    pdf-ffbe9c7b3306-p186-b001}
\proofstep
  {FABL.expect_satisfiedIndicator_maxE3Sat}
  {direct}
  {Analysis-of-Boolean-Functions-by-Ryan-ODonnell}
  {pdf-ffbe9c7b3306-p190-b002,
    pdf-ffbe9c7b3306-p190-b004}
\proofstep
  {FABL.expect_satisfiedIndicators_maxE3Sat}
  {context}
  {Analysis-of-Boolean-Functions-by-Ryan-ODonnell}
  {pdf-ffbe9c7b3306-p190-b004}
\proofstep
  {FABL.maxArity_maxCutTemplate}
  {context}
  {Analysis-of-Boolean-Functions-by-Ryan-ODonnell}
  {pdf-ffbe9c7b3306-p184-b003,
    pdf-ffbe9c7b3306-p184-b004}
\proofstep
  {FABL.maxArity_maxThreeColoringTemplate}
  {context}
  {Analysis-of-Boolean-Functions-by-Ryan-ODonnell}
  {pdf-ffbe9c7b3306-p184-b003,
    pdf-ffbe9c7b3306-p184-b004}
\proofstep
  {FABL.maxArity_maxE3LinTemplate}
  {context}
  {Analysis-of-Boolean-Functions-by-Ryan-ODonnell}
  {pdf-ffbe9c7b3306-p185-b002,
    pdf-ffbe9c7b3306-p184-b004}
\proofstep
  {FABL.maxArity_maxE3SatTemplate}
  {context}
  {Analysis-of-Boolean-Functions-by-Ryan-ODonnell}
  {pdf-ffbe9c7b3306-p184-b003,
    pdf-ffbe9c7b3306-p184-b004}
Let $V$ be a finite set, and let $P$ be a nonempty finite multiset of constraints on
variables in $V$, each using three distinct variables and a three-literal disjunction,
with each literal either the variable or its negation. Under a uniformly random
assignment of Boolean values to $V$, the expected fraction of constraints satisfied,
counting multiplicities, is $7/8$.
\end{theorem}

\begin{theorem}[Seven-eighths guarantee for Max-3SAT]
\lean{FABL.seven_eighths_le_optimum_maxE3Sat}
\proofsource{Analysis-of-Boolean-Functions-by-Ryan-ODonnell}{
  pdf-ffbe9c7b3306-p190-b002,
  pdf-ffbe9c7b3306-p190-b004
}
\label{FABL.seven_eighths_le_optimum_maxE3Sat}
\leanok
\uses{FABL.CSPInstance, FABL.CSPInstance.optimum, FABL.CSPInstance.value, FABL.CSPInstance.value_le_optimum, FABL.expect_value_maxE3Sat, FABL.maxE3SatTemplate}
\difficulty{4}
\proofstep
  {FABL.𝔽₂}
  {context}
  {Analysis-of-Boolean-Functions-by-Ryan-ODonnell}
  {pdf-ffbe9c7b3306-p015-b008}
\proofstep
  {FABL.Sign}
  {context}
  {Analysis-of-Boolean-Functions-by-Ryan-ODonnell}
  {pdf-ffbe9c7b3306-p020-b001,
    pdf-ffbe9c7b3306-p022-b004}
\proofstep
  {FABL.signValue}
  {context}
  {Analysis-of-Boolean-Functions-by-Ryan-ODonnell}
  {pdf-ffbe9c7b3306-p022-b004}
\proofstep
  {FABL.signValue_one}
  {context}
  {Analysis-of-Boolean-Functions-by-Ryan-ODonnell}
  {pdf-ffbe9c7b3306-p022-b004}
\proofstep
  {FABL.signValue_neg_one}
  {context}
  {Analysis-of-Boolean-Functions-by-Ryan-ODonnell}
  {pdf-ffbe9c7b3306-p022-b004}
\proofstep
  {FABL.monomial}
  {direct}
  {Analysis-of-Boolean-Functions-by-Ryan-ODonnell}
  {pdf-ffbe9c7b3306-p022-b001,
    pdf-ffbe9c7b3306-p022-b004}
\proofstep
  {FABL.fourierCoeff}
  {direct}
  {Analysis-of-Boolean-Functions-by-Ryan-ODonnell}
  {pdf-ffbe9c7b3306-p025-b004}
\proofstep
  {FABL.fourierSupport}
  {context}
  {Analysis-of-Boolean-Functions-by-Ryan-ODonnell}
  {pdf-ffbe9c7b3306-p022-b002}
\proofstep
  {FABL.mem_fourierSupport}
  {context}
  {Analysis-of-Boolean-Functions-by-Ryan-ODonnell}
  {pdf-ffbe9c7b3306-p022-b002}
\proofstep
  {FABL.granularityFrequencyWithin}
  {context}
  {Analysis-of-Boolean-Functions-by-Ryan-ODonnell}
  {pdf-ffbe9c7b3306-p072-b003}
\proofstep
  {FABL.mem_granularityFrequencyWithin}
  {context}
  {Analysis-of-Boolean-Functions-by-Ryan-ODonnell}
  {pdf-ffbe9c7b3306-p072-b003}
\proofstep
  {FABL.CSPTemplate}
  {direct}
  {Analysis-of-Boolean-Functions-by-Ryan-ODonnell}
  {pdf-ffbe9c7b3306-p184-b004}
\proofstep
  {FABL.CSPTemplate.maxArity}
  {direct}
  {Analysis-of-Boolean-Functions-by-Ryan-ODonnell}
  {pdf-ffbe9c7b3306-p184-b004}
\proofstep
  {FABL.CSPTemplate.unsatisfiedCost}
  {context}
  {Analysis-of-Boolean-Functions-by-Ryan-ODonnell}
  {pdf-ffbe9c7b3306-p186-b001}
\proofstep
  {FABL.CSPTemplate.toValuedCSP}
  {context}
  {Analysis-of-Boolean-Functions-by-Ryan-ODonnell}
  {pdf-ffbe9c7b3306-p186-b001}
\proofstep
  {FABL.CSPTemplate.maxArity_eq_of_forall_arity_eq}
  {context}
  {Analysis-of-Boolean-Functions-by-Ryan-ODonnell}
  {pdf-ffbe9c7b3306-p184-b004}
\proofstep
  {FABL.CSPConstraint}
  {direct}
  {Analysis-of-Boolean-Functions-by-Ryan-ODonnell}
  {pdf-ffbe9c7b3306-p185-b004}
\proofstep
  {FABL.CSPConstraint.satisfiedIndicator}
  {context}
  {Analysis-of-Boolean-Functions-by-Ryan-ODonnell}
  {pdf-ffbe9c7b3306-p186-b001}
\proofstep
  {FABL.CSPConstraint.unsatisfiedIndicator}
  {context}
  {Analysis-of-Boolean-Functions-by-Ryan-ODonnell}
  {pdf-ffbe9c7b3306-p186-b001}
\proofstep
  {FABL.CSPConstraint.toValuedTerm}
  {context}
  {Analysis-of-Boolean-Functions-by-Ryan-ODonnell}
  {pdf-ffbe9c7b3306-p186-b001}
\proofstep
  {FABL.CSPConstraint.evalSolution_toValuedTerm}
  {context}
  {Analysis-of-Boolean-Functions-by-Ryan-ODonnell}
  {pdf-ffbe9c7b3306-p186-b001}
\proofstep
  {FABL.CSPInstance}
  {direct}
  {Analysis-of-Boolean-Functions-by-Ryan-ODonnell}
  {pdf-ffbe9c7b3306-p185-b004}
\proofstep
  {FABL.CSPInstance.satisfiedCount}
  {context}
  {Analysis-of-Boolean-Functions-by-Ryan-ODonnell}
  {pdf-ffbe9c7b3306-p185-b006}
\proofstep
  {FABL.CSPInstance.value}
  {direct}
  {Analysis-of-Boolean-Functions-by-Ryan-ODonnell}
  {pdf-ffbe9c7b3306-p185-b006,
    pdf-ffbe9c7b3306-p186-b001}
\proofstep
  {FABL.CSPInstance.optimalAssignment}
  {context}
  {Analysis-of-Boolean-Functions-by-Ryan-ODonnell}
  {pdf-ffbe9c7b3306-p186-b001}
\proofstep
  {FABL.CSPInstance.optimum}
  {direct}
  {Analysis-of-Boolean-Functions-by-Ryan-ODonnell}
  {pdf-ffbe9c7b3306-p186-b001}
\proofstep
  {FABL.CSPInstance.value_le_optimum}
  {context}
  {Analysis-of-Boolean-Functions-by-Ryan-ODonnell}
  {pdf-ffbe9c7b3306-p186-b001}
\proofstep
  {FABL.maxCutTemplate}
  {direct}
  {Analysis-of-Boolean-Functions-by-Ryan-ODonnell}
  {pdf-ffbe9c7b3306-p184-b003}
\proofstep
  {FABL.maxThreeColoringTemplate}
  {direct}
  {Analysis-of-Boolean-Functions-by-Ryan-ODonnell}
  {pdf-ffbe9c7b3306-p184-b003}
\proofstep
  {FABL.maxE3LinTemplate}
  {direct}
  {Analysis-of-Boolean-Functions-by-Ryan-ODonnell}
  {pdf-ffbe9c7b3306-p184-b003,
    pdf-ffbe9c7b3306-p185-b002}
\proofstep
  {FABL.e3SatPredicate}
  {context}
  {Analysis-of-Boolean-Functions-by-Ryan-ODonnell}
  {pdf-ffbe9c7b3306-p184-b003}
\proofstep
  {FABL.e3SatNegationEquiv}
  {context}
  {Analysis-of-Boolean-Functions-by-Ryan-ODonnell}
  {pdf-ffbe9c7b3306-p190-b004}
\proofstep
  {FABL.expect_plainE3SatIndicator}
  {direct}
  {Analysis-of-Boolean-Functions-by-Ryan-ODonnell}
  {pdf-ffbe9c7b3306-p190-b002,
    pdf-ffbe9c7b3306-p190-b004}
\proofstep
  {FABL.expect_e3SatPredicate}
  {direct}
  {Analysis-of-Boolean-Functions-by-Ryan-ODonnell}
  {pdf-ffbe9c7b3306-p190-b002,
    pdf-ffbe9c7b3306-p190-b004}
\proofstep
  {FABL.maxE3SatTemplate}
  {direct}
  {Analysis-of-Boolean-Functions-by-Ryan-ODonnell}
  {pdf-ffbe9c7b3306-p184-b003,
    pdf-ffbe9c7b3306-p190-b004}
\proofstep
  {FABL.e3SatScopeSplitEquiv}
  {context}
  {Analysis-of-Boolean-Functions-by-Ryan-ODonnell}
  {pdf-ffbe9c7b3306-p185-b004}
\proofstep
  {FABL.e3SatScopeSplitEquiv_fst_apply}
  {context}
  {Analysis-of-Boolean-Functions-by-Ryan-ODonnell}
  {pdf-ffbe9c7b3306-p185-b004}
\proofstep
  {FABL.expect_satisfiedIndicator_maxE3Sat}
  {direct}
  {Analysis-of-Boolean-Functions-by-Ryan-ODonnell}
  {pdf-ffbe9c7b3306-p190-b002,
    pdf-ffbe9c7b3306-p190-b004}
\proofstep
  {FABL.expect_satisfiedIndicators_maxE3Sat}
  {context}
  {Analysis-of-Boolean-Functions-by-Ryan-ODonnell}
  {pdf-ffbe9c7b3306-p190-b004}
\proofstep
  {FABL.expect_value_maxE3Sat}
  {direct}
  {Analysis-of-Boolean-Functions-by-Ryan-ODonnell}
  {pdf-ffbe9c7b3306-p190-b002,
    pdf-ffbe9c7b3306-p190-b004}
\proofstep
  {FABL.maxArity_maxCutTemplate}
  {context}
  {Analysis-of-Boolean-Functions-by-Ryan-ODonnell}
  {pdf-ffbe9c7b3306-p184-b003,
    pdf-ffbe9c7b3306-p184-b004}
\proofstep
  {FABL.maxArity_maxThreeColoringTemplate}
  {context}
  {Analysis-of-Boolean-Functions-by-Ryan-ODonnell}
  {pdf-ffbe9c7b3306-p184-b003,
    pdf-ffbe9c7b3306-p184-b004}
\proofstep
  {FABL.maxArity_maxE3LinTemplate}
  {context}
  {Analysis-of-Boolean-Functions-by-Ryan-ODonnell}
  {pdf-ffbe9c7b3306-p184-b003,
    pdf-ffbe9c7b3306-p185-b002}
\proofstep
  {FABL.maxArity_maxE3SatTemplate}
  {context}
  {Analysis-of-Boolean-Functions-by-Ryan-ODonnell}
  {pdf-ffbe9c7b3306-p184-b003,
    pdf-ffbe9c7b3306-p190-b004}
Let $V$ be a finite set of Boolean variables, and consider a nonempty finite multiset of
clauses, each containing three literals on pairwise distinct variables, where each
literal is either a variable or its negation. Assume every variable in $V$ occurs in
some clause. The maximum, over all Boolean assignments to $V$, of the fraction of
clauses satisfied, counted with multiplicity, is at least $7/8$.
\end{theorem}

\begin{definition}[Three-literal predicate data]
\lean{FABL.MaxThreeSatPredicate}
\statementsource{Analysis-of-Boolean-Functions-by-Ryan-ODonnell}{
  pdf-ffbe9c7b3306-p184-b003,
  pdf-ffbe9c7b3306-p184-b005,
  pdf-ffbe9c7b3306-p185-b003
}
\label{FABL.MaxThreeSatPredicate}
\leanok
A predicate in this family is specified by an integer $k\in\{0,1,2\}$ and a
Boolean-valued function on the set $\{0,1,\ldots,k\}$.
\end{definition}

\begin{definition}[Evaluation of a Three-Literal Predicate]
\lean{FABL.maxThreeSatPredicateEval}
\label{FABL.maxThreeSatPredicateEval}
\leanok
\uses{FABL.MaxThreeSatPredicate}
A predicate consists of an integer $n$ with $0\le n<3$ and a Boolean pattern on $n+1$
coordinates. Given Boolean values on those coordinates, the predicate evaluates to true
exactly when at least one coordinate differs from its prescribed pattern.
\end{definition}

\begin{definition}[Maximum three-SAT template]
\lean{FABL.maxThreeSatTemplate}
\statementsource{Analysis-of-Boolean-Functions-by-Ryan-ODonnell}{
  pdf-ffbe9c7b3306-p184-b003,
  pdf-ffbe9c7b3306-p184-b005,
  pdf-ffbe9c7b3306-p185-b003
}
\label{FABL.maxThreeSatTemplate}
\leanok
\uses{FABL.CSPTemplate, FABL.MaxThreeSatPredicate, FABL.maxThreeSatPredicateEval}
The maximum three-SAT template has Boolean values and a finite family of allowed
predicates. Each predicate has arity one more than the natural number in the first
component of its label, and its evaluation is given by the corresponding maximum
three-SAT predicate evaluation.
\end{definition}

\begin{theorem}[Maximum arity of the cut template]
\lean{FABL.maxArity_maxCutTemplate}
\label{FABL.maxArity_maxCutTemplate}
\leanok
\uses{FABL.CSPTemplate.maxArity, FABL.maxCutTemplate}
\difficulty{1}
The maximum arity of the maximum-cut constraint template is $2$.
\end{theorem}

\begin{theorem}[Maximum arity of the three-coloring template]
\lean{FABL.maxArity_maxThreeColoringTemplate}
\label{FABL.maxArity_maxThreeColoringTemplate}
\leanok
\uses{FABL.CSPTemplate.maxArity, FABL.maxThreeColoringTemplate}
\difficulty{1}
The maximum arity of the maximum three-coloring constraint template is $2$.
\end{theorem}

\begin{theorem}[Maximum arity of the three-variable linear template]
\lean{FABL.maxArity_maxE3LinTemplate}
\proofsource{Analysis-of-Boolean-Functions-by-Ryan-ODonnell}{
  pdf-ffbe9c7b3306-p184-b003,
  pdf-ffbe9c7b3306-p185-b002
}
\label{FABL.maxArity_maxE3LinTemplate}
\leanok
\uses{FABL.CSPTemplate.maxArity, FABL.CSPTemplate.maxArity_eq_of_forall_arity_eq, FABL.maxE3LinTemplate, FABL.𝔽₂}
\difficulty{4}
\proofstep
  {FABL.𝔽₂}
  {context}
  {Analysis-of-Boolean-Functions-by-Ryan-ODonnell}
  {pdf-ffbe9c7b3306-p015-b008,
    pdf-ffbe9c7b3306-p148-b003}
\proofstep
  {FABL.Sign}
  {context}
  {Analysis-of-Boolean-Functions-by-Ryan-ODonnell}
  {pdf-ffbe9c7b3306-p020-b001,
    pdf-ffbe9c7b3306-p014-b002}
\proofstep
  {FABL.signValue}
  {context}
  {Analysis-of-Boolean-Functions-by-Ryan-ODonnell}
  {pdf-ffbe9c7b3306-p014-b002,
    pdf-ffbe9c7b3306-p022-b004}
\proofstep
  {FABL.signValue_one}
  {context}
  {Analysis-of-Boolean-Functions-by-Ryan-ODonnell}
  {pdf-ffbe9c7b3306-p014-b002}
\proofstep
  {FABL.signValue_neg_one}
  {context}
  {Analysis-of-Boolean-Functions-by-Ryan-ODonnell}
  {pdf-ffbe9c7b3306-p014-b002}
\proofstep
  {FABL.monomial}
  {direct}
  {Analysis-of-Boolean-Functions-by-Ryan-ODonnell}
  {pdf-ffbe9c7b3306-p022-b001,
    pdf-ffbe9c7b3306-p022-b004}
\proofstep
  {FABL.fourierCoeff}
  {direct}
  {Analysis-of-Boolean-Functions-by-Ryan-ODonnell}
  {pdf-ffbe9c7b3306-p025-b004}
\proofstep
  {FABL.fourierSupport}
  {context}
  {Analysis-of-Boolean-Functions-by-Ryan-ODonnell}
  {pdf-ffbe9c7b3306-p022-b002,
    pdf-ffbe9c7b3306-p018-b003}
\proofstep
  {FABL.mem_fourierSupport}
  {context}
  {Analysis-of-Boolean-Functions-by-Ryan-ODonnell}
  {pdf-ffbe9c7b3306-p018-b003}
\proofstep
  {FABL.granularityFrequencyWithin}
  {context}
  {Analysis-of-Boolean-Functions-by-Ryan-ODonnell}
  {pdf-ffbe9c7b3306-p072-b003}
\proofstep
  {FABL.mem_granularityFrequencyWithin}
  {context}
  {Analysis-of-Boolean-Functions-by-Ryan-ODonnell}
  {pdf-ffbe9c7b3306-p072-b003}
\proofstep
  {FABL.CSPTemplate}
  {context}
  {Analysis-of-Boolean-Functions-by-Ryan-ODonnell}
  {pdf-ffbe9c7b3306-p184-b004}
\proofstep
  {FABL.CSPTemplate.maxArity}
  {direct}
  {Analysis-of-Boolean-Functions-by-Ryan-ODonnell}
  {pdf-ffbe9c7b3306-p184-b004}
\proofstep
  {FABL.CSPTemplate.unsatisfiedCost}
  {context}
  {Analysis-of-Boolean-Functions-by-Ryan-ODonnell}
  {pdf-ffbe9c7b3306-p186-b001}
\proofstep
  {FABL.CSPTemplate.toValuedCSP}
  {context}
  {Analysis-of-Boolean-Functions-by-Ryan-ODonnell}
  {pdf-ffbe9c7b3306-p184-b004,
    pdf-ffbe9c7b3306-p186-b001}
\proofstep
  {FABL.CSPTemplate.maxArity_eq_of_forall_arity_eq}
  {context}
  {Analysis-of-Boolean-Functions-by-Ryan-ODonnell}
  {pdf-ffbe9c7b3306-p184-b004}
\proofstep
  {FABL.CSPConstraint}
  {context}
  {Analysis-of-Boolean-Functions-by-Ryan-ODonnell}
  {pdf-ffbe9c7b3306-p185-b004}
\proofstep
  {FABL.CSPConstraint.unsatisfiedIndicator}
  {context}
  {Analysis-of-Boolean-Functions-by-Ryan-ODonnell}
  {pdf-ffbe9c7b3306-p186-b001}
\proofstep
  {FABL.CSPConstraint.toValuedTerm}
  {context}
  {Analysis-of-Boolean-Functions-by-Ryan-ODonnell}
  {pdf-ffbe9c7b3306-p185-b004,
    pdf-ffbe9c7b3306-p186-b001}
\proofstep
  {FABL.CSPConstraint.evalSolution_toValuedTerm}
  {context}
  {Analysis-of-Boolean-Functions-by-Ryan-ODonnell}
  {pdf-ffbe9c7b3306-p186-b001}
\proofstep
  {FABL.maxCutTemplate}
  {context}
  {Analysis-of-Boolean-Functions-by-Ryan-ODonnell}
  {pdf-ffbe9c7b3306-p184-b003}
\proofstep
  {FABL.maxThreeColoringTemplate}
  {context}
  {Analysis-of-Boolean-Functions-by-Ryan-ODonnell}
  {pdf-ffbe9c7b3306-p184-b003}
\proofstep
  {FABL.maxE3LinTemplate}
  {context}
  {Analysis-of-Boolean-Functions-by-Ryan-ODonnell}
  {pdf-ffbe9c7b3306-p184-b003,
    pdf-ffbe9c7b3306-p185-b002}
\proofstep
  {FABL.e3SatPredicate}
  {context}
  {Analysis-of-Boolean-Functions-by-Ryan-ODonnell}
  {pdf-ffbe9c7b3306-p184-b003,
    pdf-ffbe9c7b3306-p185-b003}
\proofstep
  {FABL.maxE3SatTemplate}
  {context}
  {Analysis-of-Boolean-Functions-by-Ryan-ODonnell}
  {pdf-ffbe9c7b3306-p184-b003,
    pdf-ffbe9c7b3306-p185-b003}
\proofstep
  {FABL.maxArity_maxCutTemplate}
  {context}
  {Analysis-of-Boolean-Functions-by-Ryan-ODonnell}
  {pdf-ffbe9c7b3306-p184-b003,
    pdf-ffbe9c7b3306-p184-b004}
\proofstep
  {FABL.maxArity_maxThreeColoringTemplate}
  {context}
  {Analysis-of-Boolean-Functions-by-Ryan-ODonnell}
  {pdf-ffbe9c7b3306-p184-b003,
    pdf-ffbe9c7b3306-p184-b004}
\proofstep
  {FABL.maxArity_maxE3SatTemplate}
  {context}
  {Analysis-of-Boolean-Functions-by-Ryan-ODonnell}
  {pdf-ffbe9c7b3306-p184-b003,
    pdf-ffbe9c7b3306-p184-b004}
The maximum arity of the maximum three-variable linear constraint template is $3$.
\end{theorem}

\begin{theorem}[Maximum arity of the three-variable satisfiability template]
\lean{FABL.maxArity_maxE3SatTemplate}
\label{FABL.maxArity_maxE3SatTemplate}
\leanok
\uses{FABL.CSPTemplate.maxArity, FABL.CSPTemplate.maxArity_eq_of_forall_arity_eq, FABL.maxE3SatTemplate}
\difficulty{3}
The maximum arity of the three-variable satisfiability constraint template is $3$.
\end{theorem}
